\documentclass[10pt,twocolumn,letterpaper]{article}
\usepackage[review]{cvpr}
\usepackage{amsmath,amssymb,amsthm,booktabs,needspace}
\newtheorem{theorem}{Theorem}
\newtheorem{corollary}[theorem]{Corollary}
\newtheorem{proposition}[theorem]{Proposition}
\newtheorem{lemma}[theorem]{Lemma}
\newcommand{\R}{\mathbb R}
\newcommand{\E}{\mathbb E}
\newcommand{\Prob}{\mathbb P}
\newcommand{\norm}[1]{\lVert#1\rVert}
\newcommand{\bad}{\mathcal S_{\mathrm{bad}}}

\usepackage[pagebackref,breaklinks,colorlinks,allcolors=blue]{hyperref}
\def\confName{CVPR}\def\confYear{2027}
\renewcommand{\thesection}{S\arabic{section}}
\renewcommand{\thetheorem}{S\arabic{theorem}}
\title{Scale-Controlled Reachability and Recovery Risk\\in Regularization by Denoising\\Supplementary Derivations}
\author{Anonymous authors}
\begin{document}\maketitle
\section{Notation and Proof Interfaces}
Sections S2--S5 follow the main recovery result from model and
geometric criterion through the certified parameter box and noise bounds.
Sections S6--S7 prove the two-template and three-template scale results.
The balanced family and its platform comparisons follow in separate sections.

\section{Model and Comparison Class}\label{sec:model}
For $y=Ax^\star+\eta$, consider the unmodified update
\begin{equation}\label{eq:red}
 x_{k+1}=x_k-\gamma\{A^\top(Ax_k-y)+\lambda[x_k-D_{\sigma_k}(x_k)]\},
\end{equation}
with $x_0=A^\top y$. The data loss is $\norm{Ax-y}^2/2$, without observation-variance normalization; the denoising residual is not divided by $\sigma_k^2$. Risk means $\E\norm{\widehat x-x^\star}^2$, without dimension normalization.

The risk framework in Section~\ref{sec:criterion} applies in any
finite dimension and to any finite isotropic mixture. The explicit
geometric application below specializes it to three components in two
dimensions. The denoising prior is $\sum_{i=0}^2\pi_iN(\mu_i,\tau^2I_2)$. Writing $B=A^\top A$, $b=A^\top y$, and $v=\tau^2+\sigma^2$, exact MMSE denoising gives
\begin{align}
w_i(x,v)&=\frac{\pi_i e^{-\norm{x-\mu_i}^2/(2v)}}{\sum_j\pi_j e^{-\norm{x-\mu_j}^2/(2v)}},\\
\bar\mu_v(x)&=\sum_iw_i(x,v)\mu_i,\quad c(v)=\lambda(1-\tau^2/v),\\
T_v(x)&=x-\gamma\{Bx-b+c(v)[x-\bar\mu_v(x)]\}.\label{eq:map}
\end{align}
One update uses one denoiser call. Controllers share observations, initialization, and total budget; all comparisons use the same terminal scale. Labels appear in analysis and evaluation, never as controller inputs. A scale band is an interval of available total variances $v$. A platform holds one variance for several calls. A controller selects variances and the time of switching to the terminal scale; the observation-gated preset below is one such controller. The terminal rectangles are sufficient entry sets, not complete basins of attraction.

\paragraph{The scalar potential at a fixed scale.}
Let $p_v(x)=\sum_i\pi_i(2\pi v)^{-1}e^{-\norm{x-\mu_i}^2/(2v)}$. Since $D_\sigma(x)-x=\sigma^2\nabla\log p_v(x)$, the exact field in~\eqref{eq:map} is the gradient of
\begin{equation}\label{eq:potential}
F_v(x)=\tfrac12\norm{Ax-y}^2-\lambda\sigma^2\log p_v(x),
\qquad\sigma^2=v-\tau^2.
\end{equation}
Thus $T_v=x-\gamma\nabla F_v$ at fixed $v$, and the terminal objective is $F_{v_L}$. Varying $v$ varies this potential, including its coefficient $\sigma^2$. This identity uses the exact denoiser; it does not identify a generic RED residual with the gradient of $x^\top(x-D(x))/2$.

\begin{table}[t]
\centering\small
\setlength{\tabcolsep}{3pt}
\begin{tabular}{@{}p{.39\columnwidth}p{.56\columnwidth}@{}}
\toprule
Solver and scale access & Scope of the comparison\\
\midrule
Backprojection, restricted & Risk lower bound for arbitrary switching\\
Backprojection, with band & Finite-budget consistency via the preset\\
Least-squares start & Terminal-only succeeds in the matched diagnostic\\
Change $\lambda$ or normalize data loss & Changes the vector fields and the barrier tests\\
\bottomrule
\end{tabular}
\caption{Scope of the RED comparison. The two main policies may
use all observation information when choosing scales.}
\label{tab:solver_scope}
\end{table}
For the fixed positive-definite geometry below, least squares
has risk $\nu^2(1+\alpha^{-1})$. Observations can also identify the
component on the good event. The question is whether the
specified iterates can reach its terminal region. Table~\ref{tab:solver_scope}
separates this question from changes to the solver.

\section{Proof Interface: Barrier, Channel, and Terminal Risk}
\label{sec:criterion}
To prove the risk separation, we verify three geometric tests.
Write $\widetilde T_v$ for the implemented update, including any
allowed error. The following tests apply uniformly over observations
in an event $E$, conditional on the component under consideration.

\begin{theorem}[Barrier, channel, and terminal tests]\label{thm:criterion}
Let $s$ be a unit-direction coordinate and $Z=j$ a component of an
isotropic Gaussian mixture of width $\tau$. Assume
$s(\mu_j)=\xi+d$ with $d>0$ and the following tests hold.
\begin{enumerate}
\item A closed region $W\subseteq\{s\le\xi\}$ contains $x_0$ for
$Z=j$ and satisfies $\widetilde T_v(W)\subseteq W$ for every
restricted variance. This test is required for one component with
$\pi_j>0$.
\item For every component $z$ with $\pi_z>0$, an observation-based allocation reaches a terminal set $G_Z$
within $H$ calls on $E$. A sufficient channel test for a class
needing intermediate calls uses a strip $[s_-,s_E)\times J$
containing its initial state. Each band update either stays in
the strip or crosses into $G_Z$, and
$s^+\ge q_Cs+(1-q_C)s_*$ with $0\le q_C<1$ and $s_*>s_E$.
Then $s_*-(s_*-s_-)q_C^H\ge s_E$ suffices for entry.
For a fixed holding time, the reached part of $G_Z$ must also
be preserved by the remaining band calls.
\item For every such component, the terminal map preserves $G_Z$,
with common bounds $q,\ell,M$ below. For an analysis center $m$
and $x\in G_Z$, it satisfies
$\norm{\widetilde T_{v_L}(x)-m}\le q\norm{x-m}+\ell$ with
$0\le q<1$. The distance to $m$ at switching is at most $M$.
\end{enumerate}
Write $R_m=\E\norm{m-x^\star}^2$, $p=\Prob(E^c)$,
and let $U$ bound $\big(\E\norm{\widehat x-m}^4\big)^{1/2}$ uniformly
in the budget for the policy in tests 2--3. Then
\begin{align*}
 R_{\rm bad}&\ge\pi_j[d\sqrt{\Prob(E\mid Z=j)}-\tau]_+^2,\\
 R_{\rm band}^{1/2}&\le R_m^{1/2}+
 \left[\left(q^n M+\frac{\ell}{1-q}\right)^2+U\sqrt p\right]^{1/2}.
\end{align*}
The first bound holds for any finite budget. The second uses a
budget of $H+n$ with at least $n$ terminal calls on $E$.
\end{theorem}

\begin{proof}
Induction keeps every restricted trajectory in $W$ on $E$ for class
$j$, so its directional distance from $\mu_j$ is at least $d$ there.
The reverse triangle inequality in conditional $L^2$ subtracts at
most $\tau$ for the within-component displacement.
The affine lower comparison forces the strip exit by $H$ calls.
Iterating the terminal estimate bounds the distance to $m$ on $E$.
Cauchy--Schwarz bounds its squared contribution on $E^c$ by
$U\sqrt p$. The $L^2$ triangle inequality gives the upper bound.
\end{proof}

In a sensing eigenbasis, $1-\gamma(B_{ii}+c(v))\ge0$
makes the exact map nondecreasing in its own coordinate, since
$\partial_i\bar\mu_{v,i}=\operatorname{Var}_{w}(\mu_{Z,i})/v$.
For rectangular regions, inward face tests with an error margin
therefore suffice for the set inclusions above, including unbounded
faces. With bounded templates and a positive lower bound on $c(v)$,
the coordinate update gives a budget-uniform moment envelope.
The criterion can be used with other geometries when these tests
are verified. The structured application verifies these tests with explicit margins.

\section{Main Result: Scale Access and Recovery Risk}\label{sec:extensions}
The separation holds over a box of template geometries and masses.
We first give the aligned result. Consider three equal-width
isotropic Gaussian components with
\begin{gather}
 B=\operatorname{diag}(\alpha,1),\quad 0<\alpha\le10^{-3},
 \quad\gamma=1/4,\nonumber\\
 \mu_0=(0,0),\quad \mu_1=(1,0),\quad \mu_2=(-a,d),\nonumber\\
 a\in[.98,1.02],\quad d\in[1.98,2.02],\nonumber\\
 r=\pi_1/\pi_0\in[1.27,1.31],\quad
 h=\pi_2/\pi_0\in[1.60,1.96].\label{eq:recovery_box}
\end{gather}
Let $\tau$ be the common component width of the truth and the prior,
and let $\nu$ be the measurement noise standard deviation.
Choose $v_L,\lambda$ with
\begin{equation}\label{eq:ind_coeff}
 0<\tau^2<v_L\le10^{-3},\quad
 \lambda\le1.002,\quad c(v_L)\ge.999 .
\end{equation}
We compare policies for the same update~\eqref{eq:map} from $x_0=b$.
Both use $v\ge v_L$ and can access $b$ and their entire history.
A restricted controller can use only
\begin{equation}\label{eq:main_bad}
 [v_L,.17]\cup[2.5,\infty).
\end{equation}
The intermediate preset predicts class 2 if $b_2>1$, class 1 if
$b_2\le1$ and $b_1>\alpha/2$, and class 0 otherwise. It allocates
$0,25,0$ calls, respectively, in $[.69,.71]$, then uses $v_L$ for
the rest of the budget. The controller uses the observation alone
for this allocation.

\begin{theorem}[Risk separation under restricted scale access]\label{thm:ind_noise}
Suppose each update error is bounded by $10^{-4}$ in $\ell_\infty$.
On the event
\begin{equation}\label{eq:ind_event}
 E=\{|e_1|\le\alpha/4,\ |e_2|\le .005\},
 \qquad e=b-B\mu_Z ,
\end{equation}
every restricted controller keeps $s_k\le.3$ for all $k$ when $Z=1$.
The intermediate preset enters the correct terminal rectangle on
$E$ for each class.

Suppose $Z\sim\pi$, $u\sim N(0,\tau^2 I_2)$, and
$\eta\sim N(0,\nu^2 I_2)$ are mutually independent, with
$x^\star=\mu_Z+u$ and $y=Ax^\star+\eta$.
Write $S=(\tau^2+\nu^2)^{1/2}\le1$ and $p_E=\Prob(E^c)$. Then
\begin{align}
 p_E &\le
 2e^{-1/[32(\tau^2+\nu^2/\alpha)]}
 +2e^{-.005^2/(2S^2)},\label{eq:ind_tail}\\
 R_{\rm bad}&\ge
 \pi_1[.7\sqrt{1-p_E}-\tau]_+^2,\label{eq:ind_lower}\\
 R_{\rm mid}&\le C\big(S^2+q^{2n}+e^{-.06/v_L}
             +\delta_L^2+p_E^{1/2}\big),\qquad q=.751.\label{eq:ind_upper}
\end{align}
The lower bound is uniform over all finite budgets and restricted
controllers. The upper bound uses budget $25+n$ and terminal
error at most $\delta_L\le10^{-4}$ on the terminal rectangles defined in S5.
The constant $C$ is uniform over the stated ranges.
\end{theorem}

Fix $A$ and the mixture means and masses in~\eqref{eq:recovery_box}.
As $\tau,\nu\to0$ independently, take $v_L=S$, $\lambda=1$, and
$\delta_L=O(S)$. With $L=\log(1/S)$ and $n=\lceil L/\log(1/q)\rceil$,
\begin{align}\label{eq:fixed_geom}
 \liminf_{S\to0}\inf_{\substack{K<\infty\\
             \text{restricted controllers}}} R_K &\ge .147, &
 R_{\rm mid}&=O(S^2).
\end{align}
The infimum may choose a different controller and budget at each noise level.
All policies at that level share the prior, coefficient, and terminal variance.
The rate also holds for varying $\alpha$ if
\begin{align}
 \tau^2+\nu^2/\alpha &\le (128L)^{-1},\nonumber\\
 S^2 &\le .005^2/(8L),\label{eq:ind_regime}\\
 v_L &\le .03/L,\qquad \delta_L=O(S).\nonumber
\end{align}

\paragraph{Calibrated oracle risk.}
The known-component Gaussian conditional mean has risk
\begin{equation}\label{eq:oracle_general}
 R_{\rm oracle}=
 \frac{\tau^2\nu^2}{\alpha\tau^2+\nu^2}
 +\frac{\tau^2\nu^2}{\tau^2+\nu^2}.
\end{equation}
For fixed $\alpha>0$ this is of order $\min(\tau^2,\nu^2)$ with
constants depending on $\alpha$. The general rate in
\eqref{eq:fixed_geom} describes consistency. Under the additional
calibration $\tau=\nu=\alpha^2$, $v_L=\alpha^2+\alpha^4$ and
$\lambda=(1+\alpha^2)(1+\theta\alpha)$ for $|\theta|\le1$, exact
terminal updates yield $R_{\rm mid}\le R_{\rm oracle}+5\alpha^6$
for $\alpha\le10^{-5}$ with
$n=\lceil\log(12/\alpha^3)/\log(4/3)\rceil$. The calculation
appears in the calibrated specialization in S5.

\subsection{Proof of Theorem~\ref{thm:ind_noise}}
The coefficient condition implies $.999\le c(v)\le1.002$ for every $v\ge v_L$.
We use the margins in Section~\ref{sec:geometry_full} with the separate coordinate tolerances
in~\eqref{eq:ind_event}. At the barrier's right face, the horizontal
tolerance increases by at most $.001/4$, so $q_B>.298896-.00025/.999>.29864$.
The certified bounds $L(.17)<.287729$ and $U(2.5)<.289884$ still suffice.
The bottom barrier and channel tests have left side at least
$\gamma[(1+.999)(.01)-.005]>10^{-4}$.
For channel progress, $b_1-\alpha s\ge \alpha(1-s_E-1/4)>0$.
The comparison sequence uses this sign and the same $\kappa,\Delta_E$,
and so the 25-step bound is unchanged. The initial $s$ is between
$s_-=-.01$ and $s_E=.53$ for class 1, and $|t|\le.005$.
The certified channel top-face drift divided by $\gamma$ increases
by at most $.005$, from at most $-.01064024$.
The overshoot bound increases by at most $\gamma(.00025)$, remaining
below $.649<s_U$. The right-face test holds since
$\gamma[(\alpha+.999)(.1)-.00025]>10^{-4}$.
The classifier thresholds still distinguish the three means.
Class 0 starts within $.005$ of the origin in each coordinate, and
the class-2 initial coordinates satisfy $-.00127\le s\le .00025$
and $1.975\le t\le2.025$, inside $\widetilde G_2$.

The terminal gaps remain at least $.06$ and imply
$\Lambda_z\le4e^{-.03/v_L}\le4e^{-30}$.
Every template-to-face distance is at least $.09$, so
$.005+1.002\cdot3\cdot4e^{-30}+10^{-4}/\gamma<.999\cdot .09$.
The terminal invariance test holds with the global error budget.
In the local residual bound the linear factor is at most $1-\gamma(.999)<q$.
The entry distance to $m_Z(c_L)$ is at most 6, and
$c_L/(\alpha+c_L)\le1$. Hence on $E$ the residual norm after at least $n$
terminal calls is at most
\[
 6q^n+12e^{-.03/v_L}+(\sqrt{2}/(\gamma\cdot.999))\delta_L .
\]

Conditional on $Z$, $e=Bu+A^\top\eta$ has coordinate variances
$\alpha^2\tau^2+\alpha\nu^2$ and $\tau^2+\nu^2$ and its law is independent of $Z$.
Gaussian tails and the union bound give~\eqref{eq:ind_tail}.
On $E$ for $Z=1$, the restricted output satisfies $1-\widehat s\ge .7$.
Applying the reverse triangle inequality in conditional $L^2$ to
$(1-\widehat s+u_1)\mathbf{1}_E$ gives~\eqref{eq:ind_lower}.

For the upper bound, the exact coordinate update is a convex combination
of $x_j$ and $(b_j+c(v)\bar\mu_{v,j}(x))/(\beta_j+c(v))$, with weight
$\gamma(\beta_j+c(v))\ge .999/4$ on the latter.
The latter value has absolute value bounded by $3+2|e_j|$, using
$b_j=\beta_j\mu_{Zj}+e_j$ and the bounded template coordinates.
The same bound holds at initialization and for $m_Z(c_L)$.
With update errors, induction adds at most $10^{-4}/(\gamma\cdot.999)$
to the bound on the iterates. As $\operatorname{Var}(e_j)\le S^2\le1$,
Gaussian fourth moments give a uniform $L^4$ bound on
$\|\widehat x-m_Z(c_L)\|$, independent of $\alpha$ and of the budget.
By Cauchy--Schwarz the contribution to the squared residual on $E^c$
is at most a constant times $p_E^{1/2}$.
Finally,
\[
 \E\norm{m_Z(c_L)-x^\star}^2
   =\sum_{\beta\in\{\alpha,1\}}
     \frac{\beta\nu^2+c_L^2\tau^2}{(\beta+c_L)^2}
   \le C_0 S^2 .
\]
The squared triangle inequality gives~\eqref{eq:ind_upper}.
Under~\eqref{eq:ind_regime}, $p_E\le4S^4$, the terminal leakage is at most
$S$ before squaring, and $q^n\le S$. Since $\pi_1>.30$, the stated
separation follows. At fixed $\alpha>0$, both Gaussian tail conditions
in~\eqref{eq:ind_regime} hold for sufficiently small $S$. Also $\tau^2/v_L\le S$ for $v_L=S$, which yields the coefficient window
eventually at $\lambda=1$.


\onecolumn
\section{Computable geometry: statements and complete proofs}\label{sec:geometry_full}
\subsection{Free-ratio model}
Let
\begin{equation}\label{geo:eq:model}
 B=\operatorname{diag}(\alpha,1),\qquad
 \mu_0=(0,0),\quad\mu_1=(1,0),\quad\mu_2=(-a,d),\qquad a,d>0.
\end{equation}
Write $r=\pi_1/\pi_0$, $h=\pi_2/\pi_0$, so that
$\pi=(1,r,h)/(1+r+h)$. These are free positive inputs. The letter $h$ denotes
a prior ratio throughout this geometric criterion, never a holding time. For total denoising
variance $v=\tau^2+\sigma^2$, the exact mixture-MMSE RED map is
\begin{align}
 w_i(x,v)&=\frac{\pi_i\exp[-\norm{x-\mu_i}^2/(2v)]}
 {\sum_j\pi_j\exp[-\norm{x-\mu_j}^2/(2v)]},\qquad
 \overline\mu_v(x)=\sum_iw_i(x,v)\mu_i,\label{geo:eq:weights}\\
 T_v(x)&=x-\gamma\{Bx-b+c(v)[x-\overline\mu_v(x)]\},\qquad
 c(v)=\lambda(1-\tau^2/v).\label{geo:eq:update}
\end{align}
Initialization is $x_0=b$. Each iteration is $x_{k+1}=T_{v_k}(x_k)+e_k^{\rm upd}$.
The controller and the update perturbation may depend measurably on the
observation, history, and an independent random seed. They do not receive $Z$.
For deterministic analysis the observation event is
\begin{equation}\label{geo:eq:obs}
 b=B\mu_Z+e,\qquad \norm{e}_\infty\le\varepsilon.
\end{equation}
We first derive sufficient conditions for general $0<\alpha\le1$ and
$0<c_-\le c(v)\le c_+$, assuming
\begin{equation}\label{geo:eq:step}
 0<\gamma\le(1+c_+)^{-1},\qquad
 \norm{e_k^{\rm upd}}_\infty\le\delta.
\end{equation}
All scale intervals in this geometric criterion refer to $v$, not $\sigma$.

The numerical application uses the independent box
\begin{equation}\label{geo:eq:box}
\begin{gathered}
0<\alpha\le10^{-3},\quad a\in[.98,1.02],\quad d\in[1.98,2.02],\\
r\in[1.27,1.31],\quad h\in[1.60,1.96],\quad\theta\in[-1,1],\\
\gamma=\tfrac14,\quad \tau=\alpha^2,\quad\sigma_L=\alpha,\quad
v_L=\alpha^2+\alpha^4,\quad\lambda=(1+\alpha^2)(1+\theta\alpha).
\end{gathered}
\end{equation}
For $v\ge v_L$, $c_L=c(v_L)=1+\theta\alpha$,
$.999\le c(v)\le1.002$, and $\alpha+c(v)\ge1$.
The observation tolerance is $\varepsilon=\alpha/500$ and the global update-error
radius is $\delta=10^{-4}$. There is no balancing constraint
in \eqref{geo:eq:box}.

\subsection{Explicit geometric conditions}
For $x=(s,t)$ define the odds
\begin{align}
 R_v(s)&=r\exp[(s-\tfrac12)/v],\label{geo:eq:oddsR}\\
 H_v(s,t)&=h\exp[-(a^2+d^2+2as-2dt)/(2v)].\label{geo:eq:oddsH}
\end{align}
The horizontal posterior mean and third responsibility are
\begin{equation}\label{geo:eq:mean}
 m_v(s,t)=\frac{R_v-aH_v}{1+R_v+H_v},\qquad
 w_2(s,t,v)=\frac{H_v}{1+R_v+H_v}.
\end{equation}
These identities turn all subsequent tests into direct scalar computations.

\subsection{A common obstruction over two scale ranges}
Choose $0<\xi<1/2$, $\eta_B\ge\varepsilon$, and set
\begin{equation}\label{geo:eq:qB}
 q_B=\xi-\frac{\alpha(1-\xi)+\varepsilon+\delta/\gamma}{c_-}.
\end{equation}
Assume $0<q_B<\xi$, $\alpha+\varepsilon\le\xi$, and
\begin{equation}\label{geo:eq:bottomB}
 \gamma[(1+c_-)\eta_B-\varepsilon]\ge\delta.
\end{equation}
For positive cutoffs $V_-<V_+$ define
\begin{align}
 L(V_-)&=\frac{r\exp[-(1/2-\xi)/V_-]}
 {1+r\exp[-(1/2-\xi)/V_-]},\label{geo:eq:low}\\
 D_B&=a^2+d^2+2a\xi+2d\eta_B,\qquad
 H_B=h\exp[-D_B/(2V_+)],\label{geo:eq:highodds}\\
 U(V_+)&=\frac{r-aH_B}{1+r+H_B}.\label{geo:eq:high}
\end{align}

\begin{proposition}[Computable barrier]\label{geo:prop:barrier}
If \eqref{geo:eq:step}--\eqref{geo:eq:bottomB} hold and
\begin{equation}\label{geo:eq:barriertest}
 L(V_-)\le q_B,\qquad U(V_+)\le q_B,
\end{equation}
then, on \eqref{geo:eq:obs} with $Z=1$, every sequence of admissible scales in
$(0,V_-]\cup[V_+,\infty)$ preserves
\[
 W=\{(s,t):s\le\xi,\ t\ge-\eta_B\}
\]
from $x_0=b$, including arbitrary adaptive or randomized switching and all
perturbations satisfying \eqref{geo:eq:step}. The statement holds for every call
budget. Any minimum admissible variance is intersected with this union.
\end{proposition}

The cutoffs are explicit. If $r>q_B/(1-q_B)$, the low-scale condition is
equivalent to
\begin{equation}\label{geo:eq:lowcut}
 V_-\le\frac{1/2-\xi}{\log[r(1-q_B)/q_B]}.
\end{equation}
Otherwise it is automatic. Write
$H_{\rm req}=[r(1-q_B)-q_B]/(a+q_B)$. If $H_{\rm req}>0$ and $h>H_{\rm req}$,
the high-scale condition is equivalent to
\begin{equation}\label{geo:eq:highcut}
 V_+\ge\frac{D_B}{2\log(h/H_{\rm req})}.
\end{equation}
If $H_{\rm req}\le0$ it is automatic. When $H_{\rm req}>0$ and
$h\le H_{\rm req}$, this particular high-scale test supplies no finite cutoff.

\subsection{An intermediate progress channel and a safe switch}
Choose $s_-<1/2<s_E<1<s_U$, $\eta_C\ge\varepsilon$, $t_+\ge\varepsilon$, and a band
$I=[v_-,v_+]$. Require $\alpha-\varepsilon\ge s_-$, $\alpha+\varepsilon\le s_E$, and
$\varepsilon\le\alpha(1-s_E)$. Define
\begin{align}
 D_C&=a^2+d^2+2as_--2dt_+>0,\nonumber\\
 R_C&=r\exp[(s_--1/2)/v_-],\quad
 H_C=h\exp[-D_C/(2v_+)],\quad
 W_C=\frac{H_C}{1+R_C+H_C},\label{geo:eq:WC}\\
 \kappa&=\frac{1/4+a(a+2)W_C}{v_-}<1,\qquad\zeta=1-\kappa,\label{geo:eq:kappa}\\
 D_E&=a^2+d^2+2as_E-2dt_+>0,\nonumber\\
 R_E&=r\exp[(s_E-1/2)/v_+],\quad H_E=h\exp[-D_E/(2v_+)],\nonumber\\
 \Delta_E&=\frac{R_E-aH_E}{1+R_E+H_E}-s_E>\frac{\delta}{\gamma c_-}.
 \label{geo:eq:surplus}
\end{align}
The transverse, overshoot, and prepared-region tests are
\begin{align}
 \gamma[(1+c_-)\eta_C-\varepsilon]&\ge\delta,\label{geo:eq:Cbottom}\\
 \varepsilon-t_++\max_{c\in\{c_-,c_+\}}c(dW_C-t_+)+\delta/\gamma&\le0,\label{geo:eq:Ctop}\\
 U_E:=s_E+\gamma[(\alpha+c_+)(1-s_E)+\varepsilon]+\delta&\le s_U,\label{geo:eq:exit}\\
 \gamma[(\alpha+c_-)(s_U-1)-\varepsilon]&\ge\delta.\label{geo:eq:right}
\end{align}
Set
\begin{equation}\label{geo:eq:comparison}
 q_C=1-\gamma c_-\zeta,\qquad
 s_*=s_E+\frac{\Delta_E-\delta/(\gamma c_-)}{\zeta}.
\end{equation}
Then $0<q_C<1$ and $s_*>s_E$.

\begin{proposition}[Computable intermediate entry]\label{geo:prop:channel}
Under the preceding tests, on \eqref{geo:eq:obs} for class 1, arbitrary successive
scales in $I$ enter
\[
 P_1=[s_E,s_U]\times[-\eta_C,t_+]
\]
within $N$ calls whenever
\begin{equation}\label{geo:eq:N}
 s_*-(s_*-s_-)q_C^N\ge s_E.
\end{equation}
The first crossing cannot overshoot $P_1$, and $P_1$ remains invariant under
further scales in $I$, with the same perturbation budget. Consequently both
entrance-triggered switching and a preset of $N$ intermediate calls are valid.
\end{proposition}

The certificate may replace $\kappa$ by a larger number below one and
$\Delta_E$ by a smaller positive number still satisfying
\eqref{geo:eq:surplus}; this yields a simpler, weaker recurrence.

\subsection{A prior-mass feasibility window}
The tests expose the role of the third template without solving a trajectory.
Fix a desired responsibility cap $0<\bar w<1$ and progress margin $\Delta>0$, and
write $m_*=s_E+\Delta<1$. Define
\begin{align}
 h_B(r)&=\exp[D_B/(2V_+)]\,
 \max\left\{0,\frac{r(1-q_B)-q_B}{a+q_B}\right\},\label{geo:eq:hB}\\
 h_E(r)&=\exp[D_E/(2v_+)]\,
 \frac{(1-m_*)r\exp[(s_E-1/2)/v_+]-m_*}{a+m_*},\label{geo:eq:hE}\\
 h_C(r)&=\frac{\bar w}{1-\bar w}
 \bigl(1+r\exp[(s_--1/2)/v_-]\bigr)\exp[D_C/(2v_+)].\label{geo:eq:hC}
\end{align}
For a given geometry, the explicit nonempty interval condition
\begin{equation}\label{geo:eq:window}
 h_B(r)\le h\le\min\{h_E(r),h_C(r)\},\qquad h>0,
\end{equation}
together with
\[
 r\le\frac{q_B}{1-q_B}\exp[(1/2-\xi)/V_-]
\]
implies both barrier tests, $W_C\le\bar w$, and $\Delta_E\ge\Delta$.
The remaining channel tests then use $\bar w$ and $\Delta$ directly. If the
right side of \eqref{geo:eq:window} is negative, this sufficient construction is
infeasible. The lower mass bound makes the third component strong enough to
oppose motion at high scales; the upper mass bounds keep it sufficiently
suppressed in the intermediate channel. These are sufficient inequalities,
not a characterization of every successful scale or mixture.

\subsection{A terminal interface computed from distance gaps}
Let $G_z=\prod_{j=1}^2[\ell_{zj},u_{zj}]$ contain $\mu_z$ strictly, with
face distances
$m_{zj}=\min\{\mu_{zj}-\ell_{zj},u_{zj}-\mu_{zj}\}>0$.
For every $i\ne z$ let
\begin{equation}\label{geo:eq:gaps}
 \Delta_{zi}=\min_{x\in G_z}
 \{\norm{x-\mu_i}^2-\norm{x-\mu_z}^2\}>0,
 \quad
 \Lambda_z=\sum_{i\ne z}\frac{\pi_i}{\pi_z}
 e^{-\Delta_{zi}/(2v_L)}.
\end{equation}
The minimized expression is affine in $x$, so its minimum follows from the
signs of two coefficients and the rectangle endpoints. Let $D_\infty$ and
$D_2$ be the maximum template distances in the indicated norms, and set
$\beta_1=\alpha$, $\beta_2=1$.

\begin{proposition}[Terminal invariance and raw residual]\label{geo:prop:terminal}
Suppose $c_L>0$, $\gamma(1+c_L)\le1$, and the terminal update error is at most
$\delta_L$ in $\ell_\infty$ whenever the input state lies in $\bigcup_zG_z$.
If, for both coordinates and each $z$,
\begin{equation}\label{geo:eq:terminaltest}
 \varepsilon+c_LD_\infty\Lambda_z+\delta_L/\gamma
 \le (\beta_j+c_L)m_{zj},
\end{equation}
then $G_z$ is invariant for data satisfying \eqref{geo:eq:obs} with label $z$.
With
\[
 m_z(c_L)=(B+c_LI)^{-1}(b+c_L\mu_z),\qquad
 q_L=1-\gamma\min_j(\beta_j+c_L),
\]
an input in $G_z$ satisfying $\norm{x_0-m_z(c_L)}\le M$ obeys
\begin{equation}\label{geo:eq:terminalres}
 \norm{x_n-m_z(c_L)}
 \le q_L^nM+\frac{\gamma c_LD_2\Lambda_z+\sqrt2\delta_L}{1-q_L}.
\end{equation}
The linear center is used only in analysis; the output is the actual iterate.
\end{proposition}

\subsection{Uniform application and risk separation}
Use the terminal rectangles
\begin{equation}\label{geo:eq:rectangles}
\begin{aligned}
 G_0&=[-.1,.4]\times[-.1,.15],\\
 G_1&=[.53,1.1]\times[-.1,.15],\\
 \widetilde G_2&=[-1.12,.01]\times[1.73,2.27].
\end{aligned}
\end{equation}
Classify the observation once:
\begin{equation}\label{geo:eq:classifier}
 \widehat Z(b)=
 \begin{cases}
 2,&b_2>1,\\
 1,&b_2\le1\text{ and }b_1>\alpha/2,\\
 0,&\text{otherwise}.
 \end{cases}
\end{equation}
This is a controller input calculation, not a change to reconstruction.

\begin{theorem}[Independent-box reachability and raw-output risk]\label{geo:thm:main}
For the full box \eqref{geo:eq:box}, use the same exact denoiser, coefficient,
initialization, and terminal scale for every controller.
\begin{enumerate}
\item On \eqref{geo:eq:obs} with $\varepsilon=\alpha/500$ and $Z=1$, every controller
restricted to
\begin{equation}\label{geo:eq:newbad}
 [v_L,.17]\ \cup\ [2.5,\infty)
\end{equation}
satisfies $s_k\le.3$ for all $k$, under any adaptive switching and any
global update perturbation of size at most $10^{-4}$.
\item Compute \eqref{geo:eq:classifier} once. Allocate respectively $0,25,0$
intermediate calls to predicted classes $0,1,2$, allowing any successive
variances in $[.69,.71]$, and then use $v_L$ up to a common budget $K=25+n$.
On \eqref{geo:eq:obs}, switching occurs in the correct rectangle
\eqref{geo:eq:rectangles}; every observation receives at least $n$ terminal calls.
This remains true with global update error at most $10^{-4}$.
\item Let $A=\operatorname{diag}(\sqrt\alpha,1)$, $x^\star=\mu_Z+u$, $y=Ax^\star+\eta$,
with $Z\sim\pi$ and $Z,u,\eta$ mutually independent, where
$u,\eta\sim N(0,\alpha^4 I_2)$. Assume the terminal error is at most
$\delta_L\le10^{-4}$ locally on the union of the correct rectangles, while
the global error bound remains $10^{-4}$. For the raw output,
\begin{align}
 R-R_{\rm oracle}
 &\le\left[\sqrt{(A_n+4\sqrt2\delta_L)^2+T_\alpha}+\alpha^3\right]^2,
 \label{geo:eq:fullrisk}\\
 A_n&=6(3/4)^n+12e^{-1/(64\alpha)},\qquad
 T_\alpha=128\sqrt8\,e^{-1/(4\cdot10^6\alpha^2)},\label{geo:eq:Ar}\\
 R_{\rm oracle}&=\alpha^4\left(\frac1{1+\alpha}+\frac12\right).
 \label{geo:eq:oracle}
\end{align}
In particular, if terminal updates are exact, $0<\alpha\le10^{-5}$, and
\begin{equation}\label{geo:eq:ralpha}
 n=n_\alpha=\left\lceil\frac{\log(12/\alpha^3)}{\log(4/3)}\right\rceil,
\end{equation}
then $K=25+n_\alpha$ gives $R\le R_{\rm oracle}+5\alpha^6$.
For restricted controllers as in (i), at any finite budget (possibly depending
on $\alpha$),
\begin{equation}\label{geo:eq:badrisk}
 R_{\rm bad}\ge .30(.7-\alpha/1000)^2
 \left[1-8e^{-1/(2\cdot10^6\alpha^2)}\right]_+.
\end{equation}
Thus $\liminf_{\alpha\downarrow0}R_{\rm bad}\ge.147$, whereas the successful
protocol has oracle risk ratio tending to one.
\end{enumerate}
\end{theorem}

The perturbed risk bound also gives oracle ratio for $\delta_L=o(\alpha^2)$,
and excess risk $O(\alpha^6)$ for $\delta_L=O(\alpha^3)$ at the same budget.
The constant $5$ is asserted for exact terminal updates, not for arbitrary
nonzero $O(\alpha^3)$ errors. The global bound is needed outside the good
observation event; the sharper terminal bound need only hold locally.

\subsection{The numerical margins}
The interval calculation uses the continuous parameter box, not a grid.
Its conservative readable bounds are as follows.
\begin{center}\small
\begin{tabular}{@{}p{.48\linewidth}p{.43\linewidth}@{}}\toprule
Quantity & Certified bound or value \\\midrule
Class-1 prior mass & $\pi_1>.30023$ \\
Barrier target $q_B$, using worst-case $\alpha,\varepsilon$ & $q_B>.298896$ \\
Low/high posterior horizontal mean bounds & $L(.17)<.287729$, $U(2.5)<.289884$ \\
Uniform sufficient cutoff bounds & $V_-\le.178160$, $V_+\ge2.416038$ suffice \\
Intermediate third responsibility & $W_C<.043871<.044$ \\
Horizontal mean derivative & $\partial_s m<.558172<.56$ \\
Mean surplus at $(.53,.05)$ & $\Delta_E>.016303>.016$ \\
Comparison value after 25 calls & $s_{25}^{\rm comp}>.534113$ \\
First-crossing upper bound & $s^+\le.647953$ \\
Top-face drift, including perturbation, divided by $\gamma$ & At most $-.01064024$ \\
Minimum terminal squared-distance gap & At least $.06$ \\
Global fourth-moment envelope & $\|\norm{\widehat x-m_Z(c_L)}\|_{L^4}<8.491<9$ \\
\bottomrule\end{tabular}
\end{center}
The cutoff row reports sufficient thresholds from the tests, not sharp
boundaries of actual success or failure. The theorem uses the safer cutoffs
$.17$ and $2.5$.


\subsection{Responsibility Calculus}
For $C=\sum_iw_i(\mu_i-\bar\mu)(\mu_i-\bar\mu)^\top$,
\begin{align}
D_x\bar\mu_v&=C/v,\\
H_v&=B+c(v)(I-C/v),\\
DT_v&=I-\gamma H_v.
\end{align}
The variance derivative has the pairwise expression
\begin{align}
\partial_v\bar\mu_v={1\over2v^2}\sum_{i<j}w_iw_j(\mu_i-\mu_j)\nonumber\\
{}\times[\norm{x-\mu_i}^2-\norm{x-\mu_j}^2].
\end{align}
In particular the Jacobian is symmetric for this exact denoiser; no analogous property is assumed for an arbitrary learned denoiser. The own-coordinate map derivatives are positive in the regions used in the main text. Since $C_{12}<0$, the transformed coordinates $(s,-t)$ make the platform dynamics order preserving.

\subsection{Proofs of the geometric interfaces}
\subsection{Posterior identities and monotonicity}
For $C=\sum_iw_i(\mu_i-\overline\mu_v)(\mu_i-\overline\mu_v)^\top$, direct differentiation gives
$D\overline\mu_v=C/v$. In particular,
\begin{align}
 C_{11}&=w_1(1-w_1)+a^2w_2(1-w_2)+2aw_1w_2
 \le\tfrac14+a(a+2)w_2,\label{geo:eq:C11}\\
 C_{12}&=-dw_2[a(1-w_2)+w_1]<0.\label{geo:eq:C12}
\end{align}
Thus $m_v$ is increasing in $s$ and decreasing in $t$. The third responsibility
decreases in $s$ and increases in $t$, as also follows directly from its odds.
The scalar function $(R-aH)/(1+R+H)$ increases in $R$ and decreases in $H$.
Under \eqref{geo:eq:step}, each component of $T_v$ is nondecreasing in its own
coordinate: its derivative is
$1-\gamma(\beta_j+c(v))+\gamma c(v)C_{jj}/v\ge0$.
These statements are for each fixed scale. No differentiation of an adaptive
controller or a nonsmooth update perturbation is required.

\begin{proof}[Proof of Proposition~\ref{geo:prop:barrier}]
On $s\le\xi$, at low scales,
$m_v\le R_v/(1+R_v)\le L(V_-)$. For $s\le\xi,t\ge-\eta_B$ and
$v\ge V_+$, $R_v\le r$ and $H_v\ge H_B$, whence $m_v\le U(V_+)$.
Both bounds cover the entire unbounded region and scale intervals.
At $s=\xi$ the horizontal displacement, including its worst perturbation, is
at most
\[
 \gamma\{\alpha(1-\xi)+\varepsilon+c_-(q_B-\xi)\}+\delta=0.
\]
Own-coordinate monotonicity extends this boundary bound to every $s\le\xi$.
At $t=-\eta_B$ the vertical displacement is at least
$\gamma[(1+c_-)\eta_B-\varepsilon]-\delta\ge0$, since $dw_2\ge0$.
Again own-coordinate monotonicity extends the bound. The initial point belongs
to $W$. Induction proves the claim uniformly over all permitted decisions.
Solving the two scalar inequalities yields \eqref{geo:eq:lowcut} and
\eqref{geo:eq:highcut}.
\end{proof}

\begin{proof}[Proof of Proposition~\ref{geo:prop:channel}]
On $C=[s_-,s_E]\times[-\eta_C,t_+]$, the odds obey $R_v\ge R_C$ and
$H_v\le H_C$, so $w_2\le W_C$. To see the first inequality across $s=1/2$,
note that the lowest negative exponent occurs at $s=s_-,v=v_-$; a nonnegative
exponent is larger. The second uses $D_C>0$. Equation \eqref{geo:eq:C11} gives
$\partial_s m_v\le\kappa$. At $(s_E,t_+)$, both odds have positive exponent
gaps of the indicated signs: increasing $v$ decreases $R_v$ and increases
$H_v$. Hence $m_v(s_E,t_+)\ge s_E+\Delta_E$ uniformly across the band.
Integrating the derivative bound and using \eqref{geo:eq:C12} gives
\begin{equation}\label{geo:eq:progress}
 m_v(s,t)-s\ge\Delta_E+\zeta(s_E-s),\qquad(s,t)\in C.
\end{equation}
For class 1, $b_1-\alpha s\ge\alpha(1-s_E)-\varepsilon\ge0$. Therefore, until
first crossing,
\[
 s_{k+1}\ge s_k+\gamma c_-\{\Delta_E+\zeta(s_E-s_k)\}-\delta
 =q_Cs_k+(1-q_C)s_*.
\]
The lower horizontal face advances. The transverse bounds
\eqref{geo:eq:Cbottom}--\eqref{geo:eq:Ctop} and own-coordinate monotonicity keep
$t$ in its strip. Induction gives the comparison in \eqref{geo:eq:N}, so entry
must occur by call $N$.

At first crossing, use $m_v\le1$, $s\le s_E$, and
$1-\gamma(\alpha+c(v))\ge0$ to obtain $s^+\le U_E$. The vertical strip
is still preserved. For further intermediate calls in $P_1$, its left face
points inward by the same surplus bound. Its right face points inward by
\eqref{geo:eq:right} and $m_v\le1$. On $s\ge s_E\ge s_-$ the bound $w_2\le W_C$
on the top face remains valid, so both transverse faces are safe. This proves
invariance of $P_1$ and justifies preset switching after the first entrance.
\end{proof}

\begin{proof}[Proof of the mass window]
Solving $U(V_+)\le q_B$ for $h$ gives its lower bound in
\eqref{geo:eq:window}. Solving
$(R_E-aH_E)/(1+R_E+H_E)\ge m_*$ gives $h\le h_E(r)$.
Solving $H_C/(1+R_C+H_C)\le\bar w$ gives $h\le h_C(r)$.
Finally, solving the low-scale inequality for $r$ gives its displayed upper
bound. All transformations multiply by positive denominators.
\end{proof}

\begin{proof}[Proof of Proposition~\ref{geo:prop:terminal}]
On $G_z$, the odds imply $1-w_z\le\Lambda_z$, and consequently
$\norm{\overline\mu_v-\mu_z}_p\le D_p\Lambda_z$ for $p=2,\infty$.
Each terminal coordinate is a convex combination of its present value and
\[
 \mu_{zj}+\frac{e_j+c_L(\overline\mu_{v_L,j}-\mu_{zj})}
 {\beta_j+c_L}.
\]
The weight of this second value is $\gamma(\beta_j+c_L)$.
The endpoint slack after that convex combination absorbs the update
perturbation by \eqref{geo:eq:terminaltest}; hence $G_z$ is invariant.
Subtracting $m_z(c_L)$ from the update yields exactly
\[
 x^+-m_z(c_L)=[I-\gamma(B+c_LI)](x-m_z(c_L))
 +\gamma c_L(\overline\mu_{v_L}(x)-\mu_z)+e_k^{\rm upd}.
\]
The linear operator has norm $q_L$. Summing the resulting geometric series
proves \eqref{geo:eq:terminalres}.
\end{proof}

\subsection{Proof of the uniform application}
\subsection{Barrier and intermediate entry}
Over \eqref{geo:eq:box}, $c(v)\ge c_L\ge.999$ for $v\ge v_L$, and
$c(v)\le\lambda\le(1+10^{-6})(1+.001)<1.002$.
Also $\alpha+c(v)\ge1+(1+\theta)\alpha\ge1$.
Use $\xi=.3$, $\eta_B=.01$, $c_-=.999$, $c_+=1.002$.
In \eqref{geo:eq:qB}, replacing $\alpha,\varepsilon$ by $.001,2\cdot10^{-6}$ gives
a uniform lower bound. For the high-scale bound replace $r,a,H_B$ by
$1.31,.98,H_B^{\min}$, where
\[
 H_B^{\min}=1.60\exp\left[-\frac{1.02^2+2.02^2+.6(1.02)+.02(2.02)}{5}\right].
\]
The bounds in the table verify Proposition~\ref{geo:prop:barrier}.

For the intermediate stage use
\[
 s_-=-.01,\quad s_E=.53,\quad s_U=1.1,\quad
 \eta_C=.01,\quad t_+=.05,\quad I=[.69,.71].
\]
A uniform lower bound on $D_C$ is
$.98^2+1.98^2-2(1.02)(.01)-2(2.02)(.05)$.
A uniform lower bound on $D_E$ is
$.98^2+1.98^2+2(.98)(.53)-2(2.02)(.05)$.
Using $r\ge1.27$, $h\le1.96$, and $a\le1.02$ in the monotone odds formulas
gives the table's bounds. They allow the simpler constants
\[
 W_C\le.044,\quad\kappa\le.56,\quad\zeta=.44,\quad\Delta_E\ge.016.
\]
With the global perturbation $10^{-4}$,
\[
 q_C=.89011,\qquad s_*=.53+
 \frac{.016-10^{-4}/(.25\cdot.999)}{.44}>.565453.
\]
The comparison from $s_-=-.01$ exceeds $.534113$ at call 25.
The transverse and right-face inequalities have strict slack, and the first
crossing satisfies $s^+\le.647953$. Thus $P_1\subset G_1$ is reached and can
be maintained until the preset call 25.

\subsection{Terminal rectangles and the shorter preset}
All template-to-face distances in \eqref{geo:eq:rectangles} are at least $.1$.
Uniform squared-distance gaps are
\begin{center}
\setlength{\tabcolsep}{2pt}\small
\begin{tabular}{@{}cccccc@{}}\toprule
$\Delta_{01}$&$\Delta_{02}$&$\Delta_{10}$&$\Delta_{12}$&$\Delta_{20}$&$\Delta_{21}$\\\midrule
$.2$&$4.0908$&$.06$&$5.3856$&$1.848$&$2.828$\\\bottomrule
\end{tabular}
\end{center}
For example,
$\Delta_{20}\ge-2(1.02)(.01)+2(2.02)(1.73)-1.02^2-2.02^2=1.848$;
the dependence on $d$ decreases across this parameter interval. The other
values follow by the corresponding affine endpoint evaluations.
Every prior ratio between two components is below two and $D_2<3$, so
\begin{equation}\label{geo:eq:leaknew}
 \Lambda_z\le4e^{-.03/v_L}\le4e^{-1/(33\alpha)}.
\end{equation}
The last inequality follows from $\alpha(1+\alpha^2)\le.99$.
At $\alpha=.001$, \eqref{geo:eq:terminaltest} has ample slack even with
$\delta_L=10^{-4}$, using the conservative face distance $.09$;
the leakage bound decreases for smaller $\alpha$.

Under \eqref{geo:eq:obs}, class 0 starts in $G_0$, and class 2 starts in
$\widetilde G_2$, since
\[
 b_1=-a\alpha+e_1\in[-.001022,.000002],\qquad
 b_2=d+e_2\in[1.979998,2.020002].
\]
This proves the $0/25/0$ preset. The observation rule identifies all three
labels correctly on this event. It remains defined for every observation.

The initial distance to $m_z(c_L)$ at any of these terminal entries is below
6 (a uniform bound below 2 also follows directly). Since
$\alpha+c_L\ge1$, the linear residual factor is at most $3/4$. Furthermore,
$c_L/(\alpha+c_L)\le1$, so \eqref{geo:eq:terminalres} and
\eqref{geo:eq:leaknew} give
\[
 \norm{x_{k+n}-m_z(c_L)}\le6(3/4)^n+12e^{-1/(33\alpha)}+4\sqrt2\delta_L.
\]
Relaxing $33$ to $64$ yields \eqref{geo:eq:Ar}. If desired, uniqueness of an
exact terminal fixed point in each rectangle follows from
$\norm{C}\le D_2^2(1-w_z)<9\Lambda_z$ and
$\norm{DT_{v_L}}\le.75+\gamma c_L9\Lambda_z/v_L<.751$.
The function $e^{-.03/v}/v$ increases for $v<.03$, so its endpoint check covers
all $0<\alpha\le.001$. The risk proof only needs the residual recursion.

\subsection{Unconditional risk without substituting an estimator}
Define
\[
 E_\alpha=\{\norm{u}_\infty,\norm{\eta}_\infty\le\alpha/1000\}.
\]
Gaussian tails and independence give
$\Prob(E_\alpha^c)\le8e^{-1/(2\cdot10^6\alpha^2)}$.
Since $e=Bu+A^\top\eta$, this event implies \eqref{geo:eq:obs}.

For every observation and admissible scale the exact coordinate update is
\[
 x_j^+=(1-\gamma(\beta_j+c))x_j
 +\gamma\beta_j\frac{b_j}{\beta_j}+\gamma c\overline\mu_{v,j}.
\]
Its coefficients are nonnegative and sum to one. Starting from $b$, the
perturbed iterate therefore obeys the global envelope
\begin{equation}\label{geo:eq:envelope}
 |x_{k,j}|\le3+|u_j|+|\eta_j/A_{jj}|+4\delta,
\end{equation}
because $\gamma(\beta_j+c)\ge1/4$. The unperturbed linear center has the same
bound without $4\delta$. Minkowski's inequality gives
\[
 \big\|\norm{\widehat x-m_Z(c_L)}\big\|_{L^4}
 \le\sqrt2\{6+4\delta+2\,3^{1/4}(\alpha^2+\alpha^{3/2})\}<9.
\]
In particular its fourth moment is below $9^4<128^2$, uniformly in the call
budget and all observation-dependent decisions. Cauchy--Schwarz yields
\[
 \E[\norm{\widehat x-m_Z(c_L)}^2\mathbf1_{E_\alpha^c}]
 \le128\sqrt8\,e^{-1/(4\cdot10^6\alpha^2)}=T_\alpha.
\]
This bound uses the full observation law; it does not treat truncated or
selected noise as Gaussian.

Conditional on $(y,Z)$ and any independent controller seed, the true posterior
mean is $m_Z(1)$ and the conditional covariance is
$\alpha^4(B+I)^{-1}$. Orthogonality therefore gives exactly
\begin{equation}\label{geo:eq:orthog}
 R=R_{\rm oracle}+\E\norm{\widehat x-m_Z(1)}^2.
\end{equation}
For every fixed $c>0$, subtraction of the linear centers gives
\begin{equation}\label{geo:eq:mismatch}
 \E\norm{m_Z(c)-m_Z(1)}^2
 =\alpha^4(c-1)^2\left[
 \frac{\alpha}{(\alpha+c)^2(1+\alpha)}+
 \frac1{2(1+c)^2}\right].
\end{equation}
Here $b-B\mu_Z$ has coordinate variance $\alpha^4\beta_j(1+\beta_j)$.
Using $|c_L-1|\le\alpha$, the bracket is below one on the full box, so the
$L^2$ center mismatch is at most $\alpha^3$. Combining the good-event
residual, the global tail bound, Minkowski's inequality, and \eqref{geo:eq:orthog}
proves \eqref{geo:eq:fullrisk}.

At \eqref{geo:eq:ralpha}, $6(3/4)^n\le\alpha^3/2$. The endpoint checks at
$\alpha=10^{-5}$ show $12e^{-1/(64\alpha)}\le\alpha^3/2$ and
$2T_\alpha\le\alpha^6$. The ratios
$e^{-1/(64\alpha)}/\alpha^3$ and
$e^{-1/(4\cdot10^6\alpha^2)}/\alpha^6$ increase throughout this interval,
as differentiation of their logarithms verifies. Thus $A_n\le\alpha^3$,
and $(\sqrt{A_n^2+T_\alpha}+\alpha^3)^2\le4\alpha^6+2T_\alpha\le5\alpha^6$.
For the bad controllers, on $Z=1,E_\alpha$, the horizontal error is at least
$.7-\alpha/1000$. Since
$\pi_1\ge1.27/(1+1.27+1.96)>.30$ and $Z$ is independent of $E_\alpha$,
\eqref{geo:eq:badrisk} follows. This completes Theorem~\ref{geo:thm:main}.

\subsection{Containment of the balanced subfamily's construction}
The balanced subfamily uses
$a\in[.99,1.01]$, $d\in[1.99,2.01]$, $\omega\in[.039,.041]$,
$v_0\in[.99,1.01]$, $c_0=c(v_0)$, and
\[
 q=(1/2,c_0d\omega/(1+c_0)),\qquad
 \overline w=(1/2-(a+1)\omega,1/2+a\omega,\omega),\qquad
 \pi_i\ \propto\ \overline w_i e^{\norm{q-\mu_i}^2/(2v_0)}.
\]
Eliminating the virtual equilibrium gives
\begin{align}
 r&=\frac{1/2+a\omega}{1/2-(a+1)\omega},\label{geo:eq:oldr}\\
 h&=\frac{\omega}{1/2-(a+1)\omega}
 \exp\left[\frac{a^2+a+d^2(1-2c_0\omega/(1+c_0))}{2v_0}\right].
 \label{geo:eq:oldh}
\end{align}
The first ratio increases in $a,\omega$. The second increases in $a,d,\omega$
and decreases in $v_0,c_0$ on the original box. For the only potentially
competing derivative,
\[
 \partial_\omega\log h
 =\frac1\omega+\frac{a+1}{1/2-(a+1)\omega}
 -\frac{d^2c_0}{v_0(1+c_0)}>0.
\]
Positivity follows even after dropping its middle positive term and bounding
the last term by the original endpoints and $c_0\in[.999,1.002]$.
The exponential numerator is positive, and
$1-2c_0\omega/(1+c_0)>0$, giving the other signs. Outward-rounded endpoint
evaluation proves
\begin{align*}
 1.27514855&<r<1.29651094,\\
 1.61097793&<h<1.93714021.
\end{align*}
Thus every original mixture is contained in \eqref{geo:eq:box}. The new box is
strictly larger, for example in its $a,d$ ranges, and admits freely specified
prior ratios throughout its stated Cartesian product.

This is not a claim that the original parameterization was a measure-zero
family, nor that it had fewer independent degrees of freedom. The useful
extension is a direct model-input test, a larger certified domain and scale
union, and a proof of reachability/risk that does not depend on constructing a
balanced scale. The original settled-platform calibration result is separate:
it is not asserted for every mixture in the new box.


\paragraph{Relation to restricted contraction.}
Here the unit secant $e_1=\mu_1-\mu_0$ is attenuated by $\alpha$ under $B$.
Even on the three template points, a normalization $\kappa B$ with
restricted isometry constant $\delta<1$ must have $\kappa<2.13$ from
the secant $\mu_2-\mu_0$. Hence $\delta\ge1-2.13\alpha$.
Nearest-template projection has restricted Lipschitz constant
$\beta\ge2$ by approaching the $0/1$ midpoint from class 1's side with
class 0 as reference. Thus its sufficient condition $\delta\beta<1$
fails in this regime.

\paragraph{Certificate and arithmetic cross-check.}
The standard-library program
\path{experiments/code/scale_geometry/certificate.py}
checks 56 scalar inequalities with outward-rounded Decimal arithmetic.
The 60- and 90-digit runs reproduce all 56 inequalities.
The \texttt{evidence/scale\_geometry/} folder in the companion archive
includes both programs, expected outputs and commands. The barrier
and channel checks correspond to the geometric tests in this section,
the wider-noise margins to Section~S4, and the root brackets to
Section~S7. The JSON check names identify the corresponding bound
or root in these proofs. The separate
program \path{experiments/code/scale_geometry/certificate_crosscheck.py}
uses binary interval arithmetic in \texttt{mpmath.iv} at 70 digits
and exact rational comparisons for the terminal gaps and equalities.
It checks the same 56 inequalities, 12 margins for the wider-noise
theorem, and 20 checks of the nominal mass threshold, root brackets,
and convexity conditions. It evaluates the drift directly from
Gaussian weights for the nominal root brackets.
The cross-check recomputes scalar bounds using the analytic reductions
above. These reductions handle the continuum of scales and states
through monotonicity and bounds at extremal coordinates.
The common-platform and shared-short-platform results retain their
balanced-subfamily assumptions.

\subsection{Risk of the actual terminal iterate}
Write $c_L=c(v_L)$ and define the reference point
\[
 x_{\rm aff}=(B+c_L I)^{-1}(b+c_L\mu_Z).
\]
On the terminal rectangle, the update relative to $x_{\rm aff}$
contracts its linear part by at most $q=.751$. The remaining
term is the posterior leakage plus the local update error.
Summing the resulting geometric series after at least $n$ terminal
calls bounds $\norm{\widehat x-x_{\rm aff}}$ by
$6q^n+12e^{-.03/v_L}+4\sqrt{2}\delta_L/c_-$ on $E$.
The reference point has mean-squared error
\[
 \E\norm{x_{\rm aff}-x^\star}^2=
 \sum_{\beta\in\{\alpha,1\}}
 \frac{\beta\nu^2+c_L^2\tau^2}{(\beta+c_L)^2}=O(S^2).
\]
To bound the contribution from $E^c$, each exact coordinate update
is a convex combination of its current value and
$(b_j+c(v)\bar\mu_{v,j})/(\beta_j+c(v))$.
The latter has absolute value at most $3+2|e_j|$. Because the new
value receives weight at least $\gamma c_-$, bounded update errors
increase the uniform envelope by at most $\delta/(\gamma c_-)$.
Gaussian fourth moments and Cauchy--Schwarz then bound the squared
residual contribution on $E^c$ by $O(p_E^{1/2})$ independently of the
budget. Together these estimates yield~\eqref{eq:ind_upper} for
the raw iterate.

For restricted controllers, $1-\widehat s\ge .7$ on $E$ for class 1.
The reverse triangle inequality in conditional $L^2$ subtracts
at most $\tau$ for the within-component displacement and gives
\eqref{eq:ind_lower}. The same lower bound applies to every budget
because it comes from invariance at each step.

\twocolumn
\section{When Endpoint Scales Suffice}\label{sec:two_endpoints}
Consider instead two isotropic Gaussian components of the same width
in any dimension. Suppose their separation direction is an
eigenvector of $B$ with eigenvalue $\alpha>0$. Use the affine coordinate
$s$ along the secant, with means at $0,1$ in this coordinate, and put
$r=\pi_1/\pi_0$. Let $F_v$ denote the scalar projected update of~\eqref{eq:map}.

\Needspace{5\baselineskip}
\begin{proposition}[Endpoint dominance for two aligned templates]\label{prop:two_endpoints}
Suppose $v_L\ge\tau^2$, $v_L>0$, and $\gamma(\alpha+\lambda)\le1$.
For every $s\le1/2$, every fixed observation, and $v_H\ge v_L$,
\begin{equation}\label{eq:two_endpoints}
 \sup_{v\in[v_L,v_H]}F_v(s)=\max\{F_{v_L}(s),F_{v_H}(s)\},
\end{equation}
where $v_H=\infty$ denotes the unbounded interval and
$F_\infty=\lim_{v\to\infty}F_v$.
For $s_0\le\xi\le1/2$, the endpoints match any strict crossing
of $\xi$ within $N$ calls. When $v_H=\infty$, the upper endpoint
can be replaced by a sufficiently large finite variance above
any prescribed $V\ge v_L$.
\end{proposition}

For the proof, use inverse variance $z=1/v$. The weight of the
right template is $w(z)=(1+r^{-1}e^{Dz})^{-1}$ with $D\ge0$
before the midpoint. Put $z_H=1/v_H$ with $1/\infty=0$.
If $w(z_H)\ge s$, the product $(1-\tau^2 z)(w(z)-s)$
is largest at $z_H$, since both factors are nonincreasing where the
second is nonnegative. If $w(z_H)<s$, then $w(z)\le1/2$
throughout the interval. There $w'\le0$ and $w''\ge0$,
so the product is convex and maximized at an endpoint. The projected maps
are nondecreasing in $s$. Iterating their endpoint maximum dominates
all other trajectories up to the crossing threshold, and a finite
strict crossing persists when infinite variances are approximated
by sufficiently large finite ones.

Consequently, the half-line $s\le\xi\le1/2$ is invariant for the
given scale interval exactly when both endpoint maps are inward at $\xi$.
This statement covers strict crossing up to the midpoint in the aligned
model. The next section gives a criterion for a third template to
change the local update on that part of the state space.

\subsection{Proof of Proposition~\ref{prop:two_endpoints}}
Let $\ell=\norm{\mu_1-\mu_0}>0$, let $f=(\mu_1-\mu_0)/\ell$, and use
$s=f^\top(x-\mu_0)/\ell$. Because $Bf=\alpha f$ and the mixture
covariances are isotropic, the scalar update is closed:
\begin{align*}
 F_v(s)&=s+\gamma[b_s-\alpha s+\lambda(1-\tau^2/v)(w_v(s)-s)],\\
 b_s&=f^\top(b-B\mu_0)/\ell,\\
 w_v(s)&=\big[1+r^{-1}e^{\ell^2(1/2-s)/v}\big]^{-1}.
\end{align*}
Fix $s\le1/2$ and let $z=1/v$ and $D=\ell^2(1/2-s)\ge0$.
The only scale-dependent term is proportional to
$g(z)=(1-\tau^2 z)(w(z)-s)$ for
$z\in[z_H,1/v_L]$, where $z_H=1/v_H$ (zero for $v_H=\infty$)
and $w(z)=(1+r^{-1}e^{Dz})^{-1}$.
If $w(z_H)\ge s$, then $g(z_H)\ge0$ and $g(z)\le g(z_H)$.
While $w(z)-s$ is nonnegative both factors are nonincreasing.
Afterward the product is nonpositive.
If $w(z_H)<s$, monotonicity gives $w(z)\le w(z_H)<s\le1/2$, so
\[
 w'=-Dw(1-w)\le0,\qquad
 w''=D^2w(1-w)(1-2w)\ge0.
\]
Thus $g''=(1-\tau^2z)w''-2\tau^2 w'\ge0$. The convex function $g$
has a maximum at an endpoint. For the unbounded interval,
continuity at $z=0$ gives the same supremum. This proves
\eqref{eq:two_endpoints}.

Every $F_v$ is nondecreasing in $s$, since $\partial_s w_v\ge0$
and $1-\gamma(\alpha+c(v))\ge0$. This includes the limit $F_\infty$.
Hence the two inward boundary inequalities are necessary and sufficient
for invariance. Necessity for the limiting endpoint follows by continuity.
For the finite-budget claim, iterate the pointwise maximum of the two endpoint
maps. Until this trajectory crosses $\xi$, monotonicity and
\eqref{eq:two_endpoints} make it dominate the trajectory using arbitrary
scales, by induction. If the latter crosses, the endpoint trajectory crosses
strictly within the same budget. For finite $v_H$, all these calls are
admissible. For $v_H=\infty$, its finite sequence of endpoint choices
can be realized to arbitrary accuracy by replacing all uses of $F_\infty$
with a sufficiently large finite variance, by continuity of the finite
composition in inverse variances. This preserves the strict crossing
and proves the claim for any $V$.

\section{How Template Competition Opens a Scale Band}\label{sec:geometry}
Fix a state $(s,t)$ with $0<s<1/2$ in the three-template geometry
and take $v_L>\tau^2$.
For inverse variance $z=1/v$ and $z_L=1/v_L$, define the half squared-distance
gaps relative to template 0:
\begin{equation}\label{eq:band_gaps}
 P=\tfrac12-s,\qquad Q=\tfrac12(a^2+d^2+2as-2dt).
\end{equation}
The responsibility odds relative to template 0 are $r e^{-Pz}$ and
$h e^{-Qz}$. Write $m(z)$ for the first coordinate of the posterior
template mean $\bar\mu_v$. Increasing inverse variance from zero suppresses template 2 faster than template 1 if $Q>P$.
We give the precise threshold for that suppression to open a band.

\begin{proposition}[Exact mass window at a fixed state]\label{prop:mass_band}
Fix a posterior-mean target $\rho\in(0,1)$ and define
\begin{equation}\label{eq:mass_capacity}
 H_\rho(z)=
 \frac{(1-\rho)r\,e^{(Q-P)z}-\rho e^{Qz}}{a+\rho}.
\end{equation}
If $Q\le P$, an interior variance cannot attain $m(z)>\rho$ when
both endpoints fail to attain it. If $Q>P$, set
\begin{equation}\label{eq:band_peak}
 z_\rho=\left[\frac{1}{P}
     \log\frac{(Q-P)(1-\rho)r}{Q\rho}\right]_{[0,z_L]},
\end{equation}
where the brackets clip to the specified interval.
Both endpoints have $m\le\rho$ and some intermediate variance has
$m>\rho$ exactly when
\begin{equation}\label{eq:exact_mass_window}
 \max\{H_\rho(0),H_\rho(z_L)\}\le h < H_\rho(z_\rho).
\end{equation}
The working band is the interval determined by $H_\rho(z)>h$, mapped
back to variance by $v=1/z$.
\end{proposition}

Indeed, $m(z)>\rho$ is equivalent to
\[
 (1-\rho)r>\rho e^{Pz}+(a+\rho)h e^{-(Q-P)z}.
\]
The right side is nondecreasing if $Q\le P$ and strictly convex if
$Q>P$. For the latter, the sign of $H_\rho'$ is the sign of
$(Q-P)(1-\rho)r-Q\rho e^{Pz}$. Thus the upper mass threshold
and the peak in~\eqref{eq:band_peak} are explicit. The lower threshold
keeps the two endpoints below the target. For fixed $a$ and other
inputs, raising $h$ shrinks the working interval, while increasing
$Q$ expands it by weakening the competition at each $z>0$.
At $\rho=s$, the result exactly detects whether the prior part
of the RED update can point right when both endpoints fail to do so.
It applies to positive Gaussian component widths.

\subsection{From the mass window to an update threshold}
A desired horizontal displacement per unit step size $d_*$ gives
the normalized prior target $u=(d_*-b_1+\alpha s)/\lambda$.
Set $l=1-\tau^2 z$, $A_u=(1-s)l-u$ and
\begin{equation}\label{eq:update_ratios}
 U(z)=\frac{s l+u}{A_u},\qquad
 W(z)=\frac{(a+s)l+u}{A_u}.
\end{equation}
The exact RED displacement per unit step exceeds $d_*$ precisely when
\begin{equation}\label{eq:update_band}
 J_u(z):=U(z)e^{Pz}+hW(z)e^{-(Q-P)z}<r
\end{equation}
where $A_u>0$. This uses the full $c(v)$ and the actual data force.
For $u=0$, it recovers Proposition~\ref{prop:mass_band} at $\rho=s$.

\begin{proposition}[A convex test for the RED variance band]\label{prop:convex_band}
Let $Q>P>0$, $v_L>\tau^2$, and assume $A_u(z_L)>0$ and
$s l(z_L)+u>0$. If $u\ge0$, $J_u$ is strictly convex on $[0,z_L]$.
For $u<0$ the same conclusion holds if
\begin{align}
 P^2U_L&>2P d_1+d_2,\nonumber\\
 (Q-P)^2 W_L&>(a+1)d_2,\label{eq:band_convexity}
\end{align}
where $U_L,W_L,A_L$ denote the terminal values of $U,W,A_u$, and
\[
 d_1=|u|\tau^2/A_L^2,\qquad
 d_2=2(1-s)\tau^2 d_1/A_L.
\]
Under these conditions, a nonempty band with both endpoints failing
exists exactly when $\min J_u<r\le\min\{J_u(0),J_u(z_L)\}$.
Its endpoints are the level roots of $J_u=r$.
\end{proposition}

The minimum is obtained by clipping the zero of the increasing
derivative $J_u'$ to the inverse-variance interval, or by taking the
appropriate endpoint if no zero exists. For a test of \emph{endpoint
dominance itself}, choose $d_*$ as the larger endpoint displacement
per unit step size. Whenever the hypotheses of
Proposition~\ref{prop:convex_band} hold at that target, its test
decides whether any intermediate step exceeds both endpoints.
The convexity conditions are checked at the chosen target.

\paragraph{Additional local thresholds.}
Supplement S7 derives an explicit critical competitor mass for
improvement over the endpoint steps, and S24 treats several
competitors. The risk family below uses the full drift test
\eqref{eq:update_band}, which also applies outside that critical-mass
subfamily.

\paragraph{Safe scales for a candidate barrier.}
Consider $s\le\xi,\ t\ge-\eta_B$ with $0<\xi<1/2$.
At $(\xi,-\eta_B)$ use $d_*=-\delta/\gamma$, for coordinate
update errors bounded by $\delta$. Provided the lower face is
inward and each map is nondecreasing in its own coordinate,
the variances with $J_u\ge r$ preserve the region. The other
variances admit an outward update at its corner under the allowed
error bound. Thus the roots predict where the common barrier loses
invariance for a given observation. To certify a schedule, a
separate channel argument checks progress and transverse
invariance along the trajectory, as in Theorem~\ref{thm:ind_noise}.

At the same state, the nominal recovery-box mass $h=1.78$ and
noiseless class-1 data with $\alpha=10^{-3}$ give, at
$\lambda=1+\alpha^2$ and $\delta=10^{-4}$, a corner-outward
variance interval of approximately $(.18064,1.87551)$.
The interval comes from the full update test~\eqref{eq:update_band}.
With $\delta=0$, the positive-drift interval for the exact update
is approximately $(.18095,1.87272)$. The larger interval above
expresses failure of robust invariance. At each of its scales some
allowed perturbation points outward from the corner. Reaching that
corner and entering a correct terminal rectangle require separate
trajectory arguments.
The restricted union lies outside both intervals, and the certified
working band $[.69,.71]$ lies inside. This calculation concerns a
single model and the candidate barrier. The uniform risk theorem
uses separate bounds across the parameter box and the channel.

\subsection{Convexity and exact safe scales}
Write $l(z)=1-\tau^2 z$. The normalized prior contribution is
$l(z)(m(z)-s)$, with
\[
 m(z)=\frac{r e^{-Pz}-a h e^{-Qz}}{1+r e^{-Pz}+h e^{-Qz}}.
\]
Multiplying $l(z)(m(z)-s)>u$ by the positive responsibility
denominator gives
\[
 A_u\,r e^{-Pz}> s l+u+[(a+s)l+u]\,h e^{-Qz}.
\]
Division by $A_u>0$ proves~\eqref{eq:update_band}. In
Proposition~\ref{prop:convex_band}, $A_u$ and both ratio numerators
are positive everywhere by their terminal lower bounds.

For $u\ge0$, the ratio
$V_j(z)=(j l+u)/[(1-s)l-u]$ for $j=s$ or $a+s$ satisfies
\[
 (\log V_j)''=\tau^4
 \left[\frac{(1-s)^2}{A_u^2}-\frac{j^2}{(j l+u)^2}\right]\ge0 .
\]
Multiplication by an exponential of an affine function preserves
log convexity. For $u\tau^2>0$ the displayed inequality is strict,
so both terms are strictly convex. Otherwise $U,W$ are positive
constants and the nonzero exponents give strict convexity.

For $u<0$, the derivatives satisfy
\[
 \begin{aligned}
 U'&=-\frac{|u|\tau^2}{A_u^2},&
 U''&=-\frac{2(1-s)|u|\tau^4}{A_u^3},\\
 W'&=(a+1)U',&
 W''&=(a+1)U'' .
 \end{aligned}
\]
Since $A_u$ decreases with $z$, $|U'|\le d_1$,
$|U''|\le d_2$, $U\ge U_L$, and $W\ge W_L$.
The two exponential-weighted second derivatives are bounded below
by their exponential factors times $P^2U_L-2Pd_1-d_2$ and
$(Q-P)^2W_L-(a+1)d_2$, respectively. The second bound drops
the nonnegative term $-2(Q-P)W'$. Condition~\eqref{eq:band_convexity}
makes both strictly positive. Strict convexity and continuity give
the level interval and the minimum test, including endpoint roots.

For the candidate barrier, the posterior mean in the horizontal
coordinate decreases with $t$. Its derivative equals
\[
 \partial_t m=-\frac{d}{v}w_2[a(1-w_2)+w_1]\le0 .
\]
The horizontal step at $s=\xi$ is therefore largest at $t=-\eta_B$.
Own-coordinate monotonicity extends an inward bound on the face
to the entire half-strip, even after a perturbation of size $\delta$.
The lower-face check also extends by vertical monotonicity.
A sufficient lower-face check is
$\gamma[b_2+(1+c_-)\eta_B]\ge\delta$ with $c(v)\ge c_->0$,
because the vertical posterior mean is nonnegative. For example,
the step bound $\gamma(1+c_+)\le1$ with $0<\alpha\le1$ and $c(v)\le c_+$
gives own-coordinate monotonicity via the nonnegative diagonal
of the posterior covariance. Thus, with the lower-face check,
the horizontal inequality at the corner is necessary and sufficient
for invariance for all allowed errors. At a failing scale, choose
the corner as the state and the largest allowed positive horizontal
error to exhibit a violation.

\subsection{Critical mass for endpoint improvement}
\begin{proof}
The inequalities in~\eqref{eq:critical_regime} give a nonpositive
terminal prior contribution and a positive terminal denominator $A_u$
at the high-endpoint target. The high-endpoint level is nonnegative, so
Proposition~\ref{prop:convex_band} applies. Since $J_u(0)=r$,
strict improvement is equivalent to $J_u'(0)<0$. Differentiating
the normalized prior contribution at zero gives
$D(h)/(1+r+h)^2$, with sign opposite to $J_u'(0)$.
The polynomial starts negative and has positive quadratic
coefficient for $\tau>0$. When $\tau=0$ it is increasing affine.
At the positive root $\partial_hD>0$ and
$\partial_Q D=h[a+(1+a)r]>0$. Implicit differentiation gives
$\partial_Q h_{\rm crit}<0$.
\end{proof}

\begin{figure*}[t]
 \centering
 \includegraphics[width=.91\textwidth]{figures/variance_band.pdf}
 \caption{Two applications of the scale tests at the same state $(s,t)=(.3,-.01)$.
 (a) In the mass range $0<h\le C_s/(a+s)$, the curve $h_{\rm crit}(Q)$
 separates endpoint dominance from a strict intermediate improvement.
 The plotted gaps satisfy $Q=.8+d^2/2+.01d>.8$ for $d>0$,
 with $a=1$. The dotted line marks the nominal $Q=2.82$.
 (b) With nominal $h=1.78$, the curve is the horizontal displacement
 per unit step size with the largest allowed update error. Its positive
 interval has endpoints $.18064$ and $1.87551$ and marks failure of
 robust invariance for this candidate barrier. Blue shading is limited to the plotted portions of $[v_L,.17]\cup[2.5,\infty)$. The uniform channel analysis
 certifies the narrower green band. Both panels evaluate the formulas
 at a fixed state and use no trajectory search.}
 \label{fig:variance_band}
\end{figure*}

Figure~\ref{fig:variance_band} shows why the local and uniform
tests use different margins. At the nominal corner, a wide range
of intermediate variances can break the barrier. A single
variance band for the whole model box must also maintain the
transverse strip and sustain progress at states nearer the
terminal entrance. Those are separate checks.

\clearpage
\section{Balanced family and supplementary protocols}\label{sec:balanced_family}
\subsection{An explicit structured family}
The templates have distinct roles. Templates 0 and 1 differ along the weakly observed $s$ direction. Their observations remain distinguishable at the coupled low noise level, but the data force in the RED update is only of order $\alpha$. The third template lies to the left and above them: its responsibility contributes a negative horizontal mean and couples horizontal and vertical motion. Low scales suppress template 1 near initialization; high scales retain enough third-template competition to block rightward escape. An intermediate band reduces this competition and opens a rightward channel. The free-ratio theorem formalizes these roles for a larger parameter box.

The balance construction below has an additional role. It places a virtual equilibrium exactly at $s=1/2$, so that a small scale offset can be compared to the order-$\alpha$ data separation. This reference is principally used to analyze the separate common-platform protocol; the observable controller does not solve for or track it. The computable criteria in Section S5 remove this balancing requirement; the construction below is retained for the platform comparison.
Let $0<\alpha\le10^{-3}$ and set
\begin{align}
A&=\operatorname{diag}(\sqrt\alpha,1),&\gamma&=1/4,\\
\mu_0&=(0,0),&\mu_1&=(1,0),&\mu_2&=(-a,d),
\end{align}
where $a\in[.99,1.01]$, $d\in[1.99,2.01]$, $\omega\in[.039,.041]$, $v_0\in[.99,1.01]$, and $\theta\in[-1,1]$. Couple the parameters by
\begin{equation}\label{eq:coupling}
\tau=\rho=\nu=\alpha^2,\quad\epsilon=\alpha,\quad
\lambda=(1+\alpha^2)(1+\theta\alpha).
\end{equation}
Here $\rho$ is the true within-component standard deviation and $\nu$ the observation standard deviation. The terminal variance is $v_L=\alpha^4+\alpha^2$, and $c_L=1+\theta\alpha$.

Put $c_0=c(v_0)$ and define
\begin{align}
q&=\left(\tfrac12,\frac{c_0d\omega}{1+c_0}\right),\\
(\bar w_0,\bar w_1,\bar w_2)&=(\tfrac12-(a+1)\omega,\tfrac12+a\omega,\omega),\\
\pi_i&=\frac{\bar w_i e^{\norm{q-\mu_i}^2/(2v_0)}}{\sum_j\bar w_j e^{\norm{q-\mu_j}^2/(2v_0)}}.\label{eq:weights}
\end{align}
Then $q$ is exactly fixed for virtual data $b_*=(\alpha/2,0)$ at variance $v_0$. This is an algebraic balance, not a numerical root assumption. Uniform bounds give $\pi_i>.23$ and $\pi_1>.30$. The templates remain bounded, noncollinear, and separated, and no component mass vanishes. The weights are nevertheless designed through~\eqref{eq:weights}; this construction is a subfamily of the free-ratio box, rather than a restriction needed by the geometric criterion.

\section{Recovering the Balanced Scale from Known Weights}
Let $r=\pi_1/\pi_0$. Because $q_1=1/2$ gives equal distances to templates 0 and 1,
\begin{equation}
\omega={r-1\over2[a+r(a+1)]}.
\end{equation}
Set
\begin{align}
f(v)&={a^2+a+d^2-2d^2\omega c(v)/(1+c(v))\over2v},\\
\ell&=\log\left({\pi_2\over\pi_0}{\bar w_0\over\omega}\right).
\end{align}
The numerator of $f$ is positive and decreasing, so $f'<0$. If $f(1.01)\le\ell\le f(.99)$, there is a unique $v_0\in[.99,1.01]$ with $f(v_0)=\ell$. Known-parameter bisection finds order-$\alpha$ accuracy with logarithmically many scalar evaluations, without labels, trajectories, or fixed-point computation. This construction is used for the common-platform comparison; the coarse-band observable controller does not need it.

\section{Shared Calibrated Platform}
\begin{proposition}[A sufficient common window]\label{prop:common}
On the deterministic observation event, if every intermediate variance satisfies $|v_k-v_0|\le\alpha/1000$ for at least
\begin{equation}\label{eq:hplus}
H_+(\alpha)=\left\lceil\frac{\log(64/\alpha)}{\log(25/23)}\right\rceil,
\end{equation}
all three classes enter their correct terminal regions. With $r_\alpha$ subsequent terminal calls, the same raw-output risk bound as in the main theorem holds for $\alpha\le10^{-5}$, with total budget $H_++r_\alpha$.
\end{proposition}
\begin{proof}
On the deterministic event of the main text, suppose every intermediate step satisfies $|v_k-v_0|\le\alpha/1000$. For classes 0 and 1, the rectangle $\mathcal R=[-.1,.75]\times[-.1,.15]$ is invariant, and
\begin{equation}
H_v\succeq.32I,\quad H_v\preceq2.01I,\quad
\norm{DT_v}\le.92,\quad\norm{\partial_vg_v}<4.
\end{equation}
For the noiseless actual-class data at $v_0$, order preservation around the known virtual equilibrium gives fixed points $r_0,r_1$ with
\begin{equation}
(r_0)_1\le.5-\alpha/8,\qquad(r_1)_1\ge.5+\alpha/8.
\end{equation}
Contraction comparison, including data perturbation and varying scales, gives
\begin{equation}
\norm{x_h-r_Z}\le2(.92)^h+
{\sqrt2\alpha/500+4\alpha/1000\over.32}.
\end{equation}
The second term is below $\alpha/32$. Thus
\begin{equation}
H_+(\alpha)=\left\lceil{\log(64/\alpha)\over\log(25/23)}\right\rceil
\end{equation}
steps suffice to put the first coordinate on its correct side by at least $\alpha/16$. Class 2 enters its terminal rectangle within five steps and remains there throughout the platform. This proves a shared-platform sufficient result without requiring high-scale preheating. The terminal and risk interfaces of the main text apply unchanged after entry. At $\alpha=10^{-5}$, $H_+=188$ and $r_\alpha=129$, giving 317 calls.
\end{proof}

\subsection{Actual failure for short platforms}
For a fixed $v$ with $|v-v_0|\le\alpha/100$, define
\begin{align}
\Delta(v)&={\sqrt2\alpha/500+4|v-v_0|\over.32},\\
\Gamma(v)&={\alpha/2+\sqrt2\alpha/500+4|v-v_0|\over.32},\\
M(v)&=3\alpha/32-\Delta(v).
\end{align}
\begin{proposition}[Window-specific crossing bounds]\label{prop:cross}
For a fixed $|v-v_0|\le\alpha/100$, if $M(v)>0$, $N_{\rm cross}^+(v)$ below is sufficient for entry. If the platform has at most $N_{\rm fail}(v)$ calls below, class 1 has template error greater than $.5$ after arbitrarily many terminal calls. In the smaller window $|v-v_0|\le\alpha/1000$, these bounds both have logarithmic order in $1/\alpha$. The lower bound does not apply to the observable band controller.
\end{proposition}
\begin{proof}
When $M(v)>0$, a sufficient crossing count is
\begin{equation}
N_{\rm cross}^+(v)=\max\left\{5,\left\lceil{\log(2/M(v))\over\log(25/23)}\right\rceil\right\}.
\end{equation}
For a lower bound compare to virtual data $b_*=(\alpha/2,0)$, initialized at $b_*$ and run at $v_0$. After two steps, $s_2^*<.20$ and the transformed state remains below the virtual equilibrium. Since $\partial_sT_s>.74$,
\begin{equation}
.5-s_h^*\ge .30(.74)^{h-2},\quad h\ge2.
\end{equation}
The actual class-1 trajectory is within $\Gamma(v)$ of this reference. Hence for
\begin{equation}
h\le N_{\rm fail}(v)=\left\lfloor2+
{\log[.30/(\Gamma(v)+\alpha/100)]\over\log(1/.74)}\right\rfloor
\end{equation}
it satisfies $s_h\le.5-\alpha/100$. The terminal half-space with this upper boundary is invariant, uniformly in $t$: the correct-template responsibility at its boundary is below $.00006$, and the outward drift is negative. Arbitrarily long terminal refinement therefore still leaves template error greater than $.5$. This is a true failure result within the stated platform protocol. In its safe window, both lower and upper crossing bounds have logarithmic order; they do not apply to the observable coarse-band strategy.
\end{proof}

\section{Full-Observation Tail Control}\label{sec:balanced_tail}
For any observation, each update coordinate is a convex combination of its previous value, a data coordinate $b_j/B_{jj}$, and template coordinates, since
\begin{equation}
x_j^+=(1-\gamma(B_{jj}+c))x_j+\gamma B_{jj}{b_j\over B_{jj}}+\gamma c\bar\mu_j.
\end{equation}
All coefficients are nonnegative and sum to one. Template coordinates have absolute value at most 3. Starting at $b$, this gives
\begin{equation}
|x_{k,j}|\le3+|u_j|+|\eta_j/A_{jj}|.
\end{equation}
The same envelope bounds $m_Z(c_L)$, and the Gaussian fourth moments yield $\E\norm{\widehat x-m_Z(c_L)}^4\le128^2$. Applying Cauchy--Schwarz to the exceptional event proves the $T_\alpha$ bound in the main text without truncating unbounded observations or conditioning a selected noise distribution as though it were Gaussian.

\section{Terminal Complexity for the Observable Strategy}
For $\alpha\le10^{-5}$, on the event $\norm{u}_\infty,\norm{\eta}_\infty\le1.5\alpha^2$ (probability greater than $.5$), class 1 switches with $1-s>.35$. With $q_L=[3-(1+\theta)\alpha]/4$, the terminal recursion gives
\begin{equation}
\norm{\widehat x-x^\star}\ge[.35q_L^n-2\alpha^2]_+,
\end{equation}
where $n$ is its terminal call count. Since $n\le K$ and $\pi_1>.30$,
\begin{equation}
R\ge.15[.35q_L^K-2\alpha^2]_+^2.
\end{equation}
Thus this strategy needs $K=\Omega(\log(1/\alpha))$ to reach order-$\alpha^4$ risk. This necessity concerns terminal accuracy in this strategy, not a universal complexity lower bound for arbitrary scale-control algorithms.
\section{Balanced-Subfamily Certificate}\label{sec:certificate}
\subsection{Domain and arithmetic}
The certificates enclose the family $0<\alpha\le.001$, $a\in[.99,1.01]$, $d\in[1.99,2.01]$, $\omega\in[.039,.041]$, $v_0\in[.99,1.01]$, $\theta\in[-1,1]$. Algebraic definitions retain the prescribed coupling of weights, variance and coefficient; interval dependence may be discarded only to enlarge an enclosure. The interval implementation includes $\alpha=0$ as an enclosing endpoint in nonsingular auxiliary expressions. The original algorithm and asymptotic statements still require $\alpha>0$. Singular terminal expressions are handled analytically and by monotone endpoint estimates, not by division by an interval containing zero.

Inputs are exact decimal strings, interpreted as rationals; binary floats are rejected by the interval constructor. Addition, subtraction, multiplication and division use 50-significant-digit Decimal contexts with downward/upward rounding. Division rejects denominators containing zero. Exponential and square root use Decimal's correctly rounded nearest result and enlarge it with the adjacent representable number on each side; square-root arguments are checked nonnegative. This relies on the Python Decimal implementation. It is an outward-rounded numerical certificate, not a rational formal proof or a separately verified arithmetic library.

\subsection{Analytic reduction before interval evaluation}
No parameter grid or adaptive subdivision is used by these two certificates. Each displayed formula is evaluated over its full enclosing box. Own-coordinate derivatives are positive. The cross covariance has sign
\begin{equation}
C_{12}=-dw_2[a(1-w_2)+w_1]<0,
\end{equation}
so rectangle images are controlled by the corresponding opposite corners. The five-step third-class entrance uses successive interval rectangle images; it does not enumerate successful floating-point trajectories. Absolute perturbations $|e_j|\le2\cdot10^{-6}$ enclose the deterministic event over the whole $\alpha$ range for the corner calculations. Sharper order-$\alpha$ perturbation comparisons used in the platform proof are analytic bounds with $|e_j|\le\alpha/500$, not an inference from that absolute enclosure.

Unbounded state and scale sets are not truncated. With $r=\pi_1/\pi_0$ and $h=\pi_2/\pi_0$, on $s\le.3$, $t\ge-\alpha/500$, for $v\ge4$ the responsibility odds obey
\begin{align}
{w_1\over w_0}&\le r,\\
{w_2\over w_0}&\ge h\exp[-(a^2+d^2+.6a+2d\alpha/500)/8].
\end{align}
The mean $(R-aH)/(1+R+H)$ increases in $R$ and decreases in $H$, yielding a bound for all such $s,t,v$. For $v\le.05$, $w_1-aw_2$ is at most $r e^{(s-.5)/v}/(1+r e^{(s-.5)/v})$, bounded at $s=.3$, $v=.05$. Monotonicity of the update in its own coordinate extends the boundary control to $s\to-\infty$ and $t\to+\infty$.

Terminal estimates use positive squared-distance gaps, prior ratios below two and exponentially small leakage. The checks at $\alpha=.001$ verify inward displacement and contraction; the Gaussian-risk endpoint inequalities are checked at $\alpha=10^{-5}$. For example $e^{-1/(64\alpha)}/\alpha^3$ and $e^{-1/(4\cdot10^6\alpha^2)}/\alpha^6$ increase throughout the required interval, so their endpoint checks extend to smaller $\alpha$. The contraction correction also increases on the stated interval. These monotonicity facts are analytic interfaces, not conclusions of sampling.

\subsection{Inequalities, outputs and logical use}
Table~\ref{tab:certificate} rounds archived interval bounds conservatively for readability. The JSON files retain all digits. The structural certificate checks prior masses, common-window rectangle images and Hessian bounds, the restricted-scale barrier, terminal distance gaps, terminal contraction and risk endpoint inequalities. The calibration certificate checks the coarse-band channels, first-crossing envelopes, the small-noise probability bound and local fixed-point response. Neither file alone checks the complete analytic proof.

\begin{table*}[t]
\centering\small
\begin{tabular}{p{.23\textwidth}p{.27\textwidth}p{.38\textwidth}}
\toprule
Interface & Conservative archived bound & Role in the argument\\
\midrule
Low/high restricted scales & $m_s<.023196$ / $m_s<.168215$ & Both imply negative outward drift at $s=.3$; used with monotonicity, not a bounded state scan.\\
Class-1 intermediate channel & $w_2<.041496$, $m_s(.53,.05)-.53>.018736$ & Transverse control and $\partial_sm_s<.6$ give a uniform progress envelope.\\
First exit, class 1 & Hypothetical $s_{25}>.533605$; first crossing $s^+\le.648134$ & Entry in at most 25 calls, with overshoot below $.65$.\\
Class-2 intermediate channel & $w_2>.785419$, $m_s<-.693017$; $s_5<-.518363$ & Invariant transverse rectangle and entry in at most five calls.\\
Terminal regions & Nonbinary gap $>1.89$; contraction $<.750001$ & Exponentially concentrated responsibilities, correct stable fixed point and raw residual recursion.\\
Local common platform & $\lambda_{\min}(H)>.326902$; $-\partial_v(r_*)_1>.119760$ & Contraction and fixed-point response establish a scale-induced incorrect-side margin.\\
Wrong side after settlement & Displacement coefficient $<1.571339$; margin $>.328661\alpha$ & Connects the local response to terminal incorrect-region invariance.\\
\bottomrule
\end{tabular}
\caption{Auxiliary interval margins for the exact parameter family. The prose proof uses simpler bounds with slack. The first four rows distinguish an invariant barrier from a progress channel and controlled entrance; none asserts a phase diagram over all scales.}
\label{tab:certificate}
\end{table*}

\subsection{Reproduction and archived evidence}
The companion \texttt{evidence/} directory contains \texttt{baseline/code/interval.py}, \texttt{baseline/code/structural\_certificate.py}, and \texttt{code/calibration\_certificate.py}. From that directory, with Python 3.12 and assertions enabled, run
\begin{flushleft}\scriptsize\ttfamily
python baseline/code/structural\_certificate.py\\
python code/calibration\_certificate.py
\end{flushleft}
Only the standard library is required. Do not run with \texttt{-O}: failed inequalities must stop execution. The structural JSON is written to \texttt{baseline/results/structural\_certificate.json}; the calibration JSON to \texttt{results/calibration\_certificate.json}. Both archived reruns have status \texttt{PASS}. Reproduction should use a fresh extracted copy so that archived evidence is preserved. The accompanying README gives working-directory and output details. Figure and risk diagnostics are generated by a separate NumPy/Matplotlib script with fixed parameters and fresh Gaussian samples; their successful execution is not part of the continuous-domain certificate.

\section{Settled-Platform Trapping Details}\label{sec:platform_proof}
For the virtual data $b_*=(\alpha/2,0)$, the platform map preserves $\mathcal R$ and is $.92$-Lipschitz. The implicit-function estimate and $\norm{\partial_vg_v}<4$ give
\begin{equation}
\norm{r_*(v)-q}\le12.5|v-v_0|\le.0125.
\end{equation}
Thus $r_*(v)$ lies in $[.48,.52]\times[.02,.06]$. Over this smaller box, the calibration certificate verifies $\partial_vg_1>.12$, $\partial_vg_2<-.16$. Since $H_{12}>0$ and $H$ is positive definite,
\begin{equation}
-\partial_v(r_*)_1\ge{\partial_vg_1\over H_{11}}
\ge{.12\over\alpha+c}>.1.
\end{equation}
The actual fixed points differ by at most
\begin{equation}
\norm{r_Z(v,b)-r_*(v)}\le C_b\alpha,\qquad
C_b={1/2+\sqrt2/500\over.32}<1.572.
\end{equation}
After $N_{\rm settle}$ steps the distance to the actual fixed point is at most $\alpha/10$. At $v-v_0\ge20\alpha$, class 1 therefore has $s<.5-.328\alpha$, inside the terminal half-space $s\le.5-\alpha/100$. This half-space is invariant for all $t$. At its boundary $w_1<.00006$, and the outward drift, divided by $\gamma$, is at most
\begin{equation}
\alpha(.502+\alpha/100)+.999(.00006-.49999)<0.
\end{equation}
The first-coordinate own derivative is positive, so this bound extends to the entire half-space.

At $v_0-v\ge20\alpha$, class 0 has $s>.5+.328\alpha$. Its state remains in $\mathcal R$ at switching and lies in
\begin{equation}
W_0^{\rm wrong}=[.5+\alpha/100,1.1]\times[-.1,.15].
\end{equation}
In this rectangle,
\begin{align}
{w_0\over w_1}&\le e^{-1/[100\alpha(1+\alpha^2)]},\\
{w_2\over w_1}&\le2e^{-5.32/[2\alpha^2(1+\alpha^2)]},
\end{align}
so $m_s>.9999$. At the left boundary the rightward drift is at least
\begin{equation}
-\alpha(.502+\alpha/100)+.999(.9999-.50001)>0.
\end{equation}
The other boundaries are preserved by coordinate convexity and the exponential bound on $dw_2$. Thus $W_0^{\rm wrong}$ is invariant. In either case the first-coordinate distance to the true template remains greater than $.5$, even with unlimited terminal calls. For a continuous signal on $E_\alpha$, the same conclusion holds because its displacement from the template is at most $\alpha/1000$, smaller than the separator margin.

\section{Preset and Perturbed Dynamics}\label{sec:stability_proof}
This section uses the interval inequalities of Section~\ref{sec:certificate} and elementary comparisons. Put $d_*=10^{-4}$, $\gamma=1/4$, and write an implemented step as $T_v(x)+d_k$. All admissible scales satisfy $.999\le c(v)\le1.002$ and $\alpha+c(v)\ge1$. The exact map has positive own-coordinate derivatives. Face bounds for the exact map therefore extend to each corresponding coordinate extremum before the signed error $\pm d_*$ is applied. This argument uses values of the perturbed update and requires no regularity of its derivatives.

\paragraph{Unbounded wrong region.}
Use $s\le.3,\ t\ge-.01$, which contains the class-1 initialization. The low-variance responsibility bound on $m_s$ is independent of $t$. At high variance, extend the original bound from $t\ge-\alpha/500$ down to $-.01$. Since $|\partial_t m_s|=|C_{12}|/v\le9/4$ for $v\ge4$, the new upper bound is $.168215+(.01)(9/4)<.20$. At the $s=.3$ face, the outward step is bounded by $\gamma(.702\alpha-.1c)+d_*<0$. At the lower face $t=-.01$, the inward step is at least $\gamma((1+c).01-2\cdot10^{-6})-d_*>0$. Thus arbitrary restricted switching preserves this larger region. The same class-1 risk lower bound follows because its first-coordinate constraint is unchanged.

\paragraph{Class 1 before and after first crossing.}
In the channel of the main proof, at the upper face $t=.05$ the drift before multiplication by $\gamma$ is at most
\[
2\cdot10^{-6}-.05+c(2.01\cdot.041496-.05)<-.016.
\]
The lower face has the inward margin given above. The bound $m_s-s\ge.23-.4s$ up to $s=.53$ gives an inward left face and the comparison
\[
s_{k+1}\ge .9001s_k+.0574425-d_* .
\]
After 25 steps the certified unperturbed comparison lower bound $.533605$ loses at most $d_* /(1-.9001)<.001002$, so the first crossing occurs by step 25. At the crossing the upper envelope is $.648134+d_*<.65$. The rectangle $[.53,1.1]\times[-.01,.05]$ is invariant at intermediate scales. At $s=.53$ use $m_s(.53,.05)-.53>.018$ and $b_1-\alpha s>0$. At $s=1.1$ use $m_s\le1$ and $b_1-\alpha s<0$. The same transverse bounds apply because $w_2$ decreases in $s$. Each inward coordinate margin exceeds $d_*$. Hence the preset switch stays inside $G_1$ after first crossing.

\paragraph{Class 2 before and after first crossing.}
On $[-1.1,.01]\times[1.75,2.1]$, use $w_2>.78$, $-a\le m_s<-.68$ and $b=B\mu_2+e$. At $t=1.75$ the drift is greater than $1.99-2\cdot10^{-6}-1.75+1.002(1.99\cdot.78-1.75)>.04$. At $t=2.1$ use $b_2\le2.010002$ and $d w_2\le2.01$. At the left face use $m_s\ge-1.01$, and at either right face $s=.01$ or $s=-.5$ use $m_s<-.68$ and $b_1-\alpha s<0$. These margins exceed $d_*$ after multiplication by $\gamma$. Until crossing $s=-.5$, the upper comparison loses at most $d_* /(1-.75025)<.000401$ relative to the exact-step envelope, so step 5 suffices. The narrower rectangle with right face $-.5$ keeps the preset inside $G_2$. Class 0 starts in $G_0$ and has zero intermediate calls. The same reasoning with $d_*=0$ proves the exact-step preset.

\paragraph{Terminal invariance and residual.}
Work on the fixed boxes $G_Z\subset Q_Z$. The exact convex-combination target of each terminal coordinate is within $3\cdot10^{-6}$ of its template coordinate, by the observation bound and the leakage estimate in the main text. Each template is at least $.09$ from the faces of its box. The damping coefficient $\gamma(B_{jj}+c_L)\ge1/4$ then gives an inward face margin of at least $(.09-3\cdot10^{-6})/4>d_*$. The perturbed iterate remains in $G_Z$. The affine error recursion relative to $m_Z(c_L)$ has linear norm at most $3/4$ and added error norm at most $\sqrt2\delta_L$. Summing this error gives $4\sqrt2\delta_L$ in addition to $A_r$. This bound controls the terminal output even for perturbations that vary across calls.

\paragraph{Exceptional observations.}
For all observations, the convex-coordinate envelope in Section~\ref{sec:balanced_tail} enlarges by at most $4d_*$ since the damping is at least $1/4$ at each step. For $\alpha\le10^{-3}$ the scalar Gaussian $L^4$ norms of $u_j$ and $\eta_j/A_{jj}$ are at most $3^{1/4}\alpha^2$ and $3^{1/4}\alpha^{3/2}$. Thus
\[
\big\|\norm{\widehat x-m_Z(c_L)}\big\|_{L^4}
\le\sqrt2[6+4d_*+2\,3^{1/4}(\alpha^2+\alpha^{3/2})]<9.
\]
This suffices for the original $T_\alpha$ by Cauchy--Schwarz. Conditional-mean orthogonality and the coefficient mismatch bound now give the perturbed risk inequality. For the exact preset the same argument uses $d_*=\delta_L=0$ and recovers the original budget and risk constant.

\paragraph{Log-weight perturbations.}
Let each denoiser log weight change by $\ell_i$ with $|\ell_i|\le10^{-4}$. Differentiating the posterior mean along the path of log tilts gives $\sum_i w_i(\mu_i-\bar\mu)\ell_i$. The diameter bound $3$ yields an update error at most $3\gamma c(v)10^{-4}<d_*$ globally. In $G_Z$ the tilted odds for each competitor increase by at most $\exp(2\cdot10^{-4})$. Both means are exponentially close to $\mu_Z$ by the terminal leakage bound, so the update error there is at most $7e^{-1/(33\alpha)}$. Use this sharper $\delta_L$ in the terminal recursion and the global $d_*$ on the exceptional event. This permits a fixed log-weight radius even for the asymptotic oracle ratio.

\section{A Local Obstruction to Shared Short Platforms}
\begin{proposition}\label{prop:shared_short}
Fix $v=.70$ and $\alpha\in[9,11]\cdot10^{-6}$. Allow each of $a,d,\omega,v_0,\theta$ to vary within $10^{-5}$ of its nominal value $(1,2,.04,1,0)$. For each shared $h\in\{0,\ldots,25\}$, the exact RED iterates on the stated deterministic observation event have template error at least $.503$ for class 1 if $h\le12$, or class 0 if $h\ge13$, after any number of terminal calls.
\end{proposition}
\begin{proof}
The shared-switch certificate propagates the backprojection rectangle for each class through 25 calls, using the two opposite corner evaluations of Section~\ref{sec:certificate}. It encloses both observation errors by $[-2.2,2.2]\cdot10^{-8}$. Each switching state is in $[-.1,1.1]\times[-.1,.15]$. For class 1 at $h\le12$ it verifies $s<.497$. For class 0 at $h\ge13$ it verifies $s>.503$.

The resulting terminal rectangles are $[-.1,.497]\times[-.1,.15]$ and $[.503,1.1]\times[-.1,.15]$. Each is contained in the corresponding $Q_j$ for the incorrectly selected template $j$. The leakage bound uses only the current state and therefore holds for these observations. Relative to the incorrect template, $b-B\mu_j$ has coordinate error at most $\alpha+\alpha/500$. The exact convex-combination target then differs from that template by less than $.002$, while the template is at least $.1$ from the rectangle faces. Both rectangles remain invariant. Their horizontal boundaries give the stated error. The certificate uses the Decimal kernel of Section~\ref{sec:certificate}. From the repository root, run \texttt{experiments/code/shared\_switch\_certificate.py} with \texttt{--out} pointing to a fresh JSON path. It encloses a continuous parameter box at each integer $h$.
\end{proof}

\begin{figure*}[t]

\centering
\includegraphics[width=\textwidth]{figures/mechanism.pdf}
\caption{Illustration on the balanced subfamily. Geometry, drift and raw RED trajectories at nominal geometry and $\alpha=10^{-3}$, with noiseless template observations. (a) Templates (stars), backprojections (circles), entrances and the clipped class-1 incorrect region. (b) A drift slice at $t=.05$: low/high scales point inward at $s=.3$; $v=.70$ advances through the channel. (c) Matched 80-call trajectories: low/high policies use $v=.05/4$ for 25 calls then terminal refinement; the earlier entrance-feedback policy switches from $v=.70$ on entrance ($h=16$, circle). Diagnostics illustrate the mechanism; uniform coverage comes from the proof.}
\label{fig:mechanism}
\end{figure*}

\section{Balanced-Subfamily Protocols}
This section and the following platform comparisons retain the balanced family of Section~\ref{sec:balanced_family}, the original entrance $G_2$, and the narrower restricted union. They do not assert platform calibration or shared-short-platform failure on the full independent box. The free-ratio theorem and its shorter preset appear in S5.

We first work on the deterministic observation event
\begin{equation}\label{eq:event}
b=B\mu_Z+e,\qquad\norm{e}_\infty\le\alpha/500.
\end{equation}
Define the available restricted scales by
\begin{equation}\label{eq:bad}
\bad=\{\sigma\ge\epsilon:v\le.05\}\cup\{\sigma:v\ge4\}.
\end{equation}
The terminal scale belongs to this set.

\subsection{An observable controller}
Compute the target-region index
\begin{equation}\label{eq:classifier}
\widehat Z(b)=\begin{cases}
2&b_2>1,\\1&b_2\le1,\ b_1>\alpha/2,\\0&\text{otherwise}.
\end{cases}
\end{equation}
This is an observation rule, not access to $Z$. Use the entrance sets
\begin{align}
G_0&=[-.1,.4]\times[-.1,.15],\\
G_1&=[.53,1.1]\times[-.1,.15],\\
G_2&=[-1.1,-.5]\times[1.5,2.1].
\end{align}
At each iteration before switching, check whether $x_k\in G_{\widehat Z}$. If so, switch permanently to $\sigma=\epsilon$. Otherwise use any $v_k\in[.69,.71]$ and check again. Force terminal switching after 25 updates, including on exceptional observations. Stop after exactly $K=25+r$ total updates. Thus early entrance yields additional terminal refinement. A deterministic implementation may simply take $v_k=.70$.

\begin{theorem}[Scale availability and raw-output risk]\label{thm:main}
Throughout the balanced family of Section~\ref{sec:balanced_family}, the following hold.
\begin{enumerate}
\item On~\eqref{eq:event}, for $Z=1$, every controller using only $\bad$ satisfies $x_{k,1}\le.3$ for every $k\ge0$. Scales may depend on observations and history.
\item On~\eqref{eq:event}, the observable controller has $\widehat Z=Z$ and enters the correct terminal invariant region after respectively $0$, at most $25$, and at most $5$ intermediate calls for $Z=0,1,2$.
\item Suppose $Z,u,\eta$ are mutually independent, with $Z\sim\pi$ and $u,\eta\sim N(0,\alpha^4I_2)$. Set $x^\star=\mu_Z+u$. Let
\begin{equation}
r_\alpha=\left\lceil\frac{\log(12/\alpha^3)}{\log(4/3)}\right\rceil,\quad K_\alpha=25+r_\alpha.
\end{equation}
For $0<\alpha\le10^{-5}$, the raw output obeys
\begin{align}
R_{\rm coarse}&\le R_{\rm oracle}+5\alpha^6,\label{eq:risk}\\
R_{\rm oracle}&=\alpha^4\left(\frac1{1+\alpha}+\frac12\right),
\end{align}
and $R_{\rm coarse}/R_{\rm oracle}\to1$. Every restricted controller has $\liminf_{\alpha\to0}R_{\rm bad}\ge.147$.
\end{enumerate}
\end{theorem}

The oracle knows the true component and returns its conditional mean; it is a risk benchmark only. Least squares has risk $\alpha^3+\alpha^4$. Hence the restricted algorithm's failure is compatible with informative observations. The controller deliberately uses this information to choose when to stop the intermediate stage, while all reconstruction steps remain~\eqref{eq:red}. No additional linear reconstruction is substituted for its output.


\begin{table*}[t]
\centering\small
\setlength{\tabcolsep}{4pt}
\begin{tabular}{p{.16\textwidth}p{.24\textwidth}p{.22\textwidth}p{.29\textwidth}}
\toprule
Protocol & Available intermediate scales & Switching rule & Current guarantee\\
\midrule
Restricted controller & $v\le.05$ with $\sigma\ge\epsilon$, or $v\ge4$ & Arbitrary observation- and history-dependent control & Class 1 remains in $\mathcal W_\alpha$ for all iterations; positive limiting risk (Theorem~\ref{thm:main}).\\[3pt]
Intermediate band & $v\in[.69,.71]$, then $v_L$ & Observable target and entrance detection; 25-step cap & At most 25 calls to enter; logarithmic refinement gives asymptotically oracle raw-output risk (Theorem~\ref{thm:main}).\\[3pt]
Settled common platform & One shared $|v-v_0|\le10^{-3}$, then $v_L$ & Hold at least $N_{\rm settle}$ before switching & Locally requires $|v-v_0|<20\alpha$ (Proposition~\ref{prop:calibration}); a narrower $\alpha/1000$ window is sufficient after $H_+$ calls (Proposition~\ref{prop:common}).\\
\bottomrule
\end{tabular}
\caption{Protocols share the RED update, initialization, coefficient and terminal scale. Theorems use the stated deterministic observation event; unconditional risk additionally uses the coupled Gaussian model. The common-platform necessary and sufficient statements have their own holding-time conditions.}
\label{tab:protocols}
\end{table*}

\section{Proof Mechanism}
The proof has four interfaces: a common incorrect invariant region, a uniformly advancing intermediate channel, posterior concentration in terminal regions, and a conditional-mean risk decomposition. Auxiliary inequalities are verified over the full parameter box by the interval procedures described in Supplementary Section~\ref{sec:certificate}.

\subsection{A common incorrect region}
\begin{lemma}[A common boundary barrier]\label{lem:barrier}
For class 1 on~\eqref{eq:event}, let $\mathcal W_\alpha=\{s\le.3,\ t\ge-\alpha/500\}$. Every scale in $\bad$ preserves this set. In particular, any observation- and history-dependent switching within $\bad$ preserves it from $x_0=b$.
\end{lemma}
\begin{proof}
Write $x=(s,t)$ and $m_s=w_1-aw_2$. For $Z=1$,
\begin{align}
s^+&=[1-\gamma(\alpha+c)]s+\gamma(b_1+cm_s),\\
t^+&=[1-\gamma(1+c)]t+\gamma(b_2+cdw_2).
\end{align}
The coordinate coefficients are nonnegative. On
$\mathcal W_\alpha$, low scales give $m_s<.025$ and high scales give $m_s<.20$. These are upper bounds on the horizontal template mean; they control the sign of the boundary drift. At $s=.3$ the outward drift is at most
\begin{equation}
\gamma\{.702\alpha+c(.20-.3)\}<0.
\end{equation}
Since $dw_2\ge0$, the lower $t$ boundary is preserved as well. Positivity of the own-coordinate derivatives extends these boundary bounds to the unbounded region. Initialization lies in $\mathcal W_\alpha$, so every restricted update preserves it, regardless of how a controller selects scales. The lower bound $\sigma\ge\epsilon$ is essential to the coefficient bounds.
\end{proof}

\subsection{Intermediate progress and entry}
\begin{lemma}[A progress channel with controlled exit]\label{lem:entry}
On~\eqref{eq:event}, the observable target equals $Z$. Every successive choice of $v\in[.69,.71]$ keeps class 1 inside a transverse channel until its first crossing of $s=.53$. This crossing occurs within 25 steps, with $s<.65$, and enters $G_1$. Class 2 enters $G_2$ within five steps, while class 0 starts in $G_0$.
\end{lemma}
\begin{proof}
Equation~\eqref{eq:event} immediately implies $\widehat Z=Z$. For the difficult class 1, before reaching $s=.53$, consider
$[-.01,.53]\times[-.01,.05]$. Responsibility monotonicity gives
\begin{equation}
w_2<.05,\quad m_s(.53,.05)-.53>.018.
\end{equation}
With $C=\sum_iw_i(\mu_i-\bar\mu)(\mu_i-\bar\mu)^\top$, we have $\partial_sm_s=C_{11}/v<.6$ and $\partial_tm_s=C_{12}/v<0$. Thus
\begin{equation}
m_s-s\ge .018+.4(.53-s)=.23-.4s.
\end{equation}
The transverse boundaries point inward, and $c\ge.999$. Therefore, until entrance, a scalar comparison with $q_C=1-\gamma(.999)(.4)=.9001$ and $\beta=.23/.4=.575$ gives
\begin{align}
s_{k+1}&\ge q_Cs_k+(1-q_C)\beta,\\
s_h&\ge\beta-(\beta+2\cdot10^{-6})q_C^h.
\end{align}
The last expression exceeds $.53$ at $h=25$. The one-step upper envelope gives $s^+<.65$ at the first crossing, so the iterate enters $G_1$ without overshooting its terminal region. The mechanism needs both progress and an exit bound; positive drift alone would not establish entrance.

For class 2 use the invariant rectangle $[-1.1,.01]\times[1.75,2.1]$. At its worst responsibility corner, $w_2>.78$ and $m_s<-.68$. Before reaching $s\le-.5$,
\begin{equation}
s_h\le-.68+(.68+.000002)(.75025)^h.
\end{equation}
At $h=5$ this is below $-.5$, yielding entry into $G_2$. Class 0 starts in $G_0$. These arguments hold uniformly for any successive variances in the entire band, not merely a sampled collection of scales. Supplementary Section~\ref{sec:certificate} records the numerical margins separately from these geometric conditions.
\end{proof}

\subsection{Terminal regions and refinement}
\begin{lemma}[Distance gaps imply raw-output refinement]\label{lem:terminal}
For the sets
\begin{align}
Q_0&=[-.1,.5-\alpha/32]\times[-.1,.15],\\
Q_1&=[.5+\alpha/32,1.1]\times[-.1,.15],\quad Q_2=G_2.
\end{align}
each $G_Z$ lies in $Q_Z$, and the terminal map preserves $Q_Z$ and has a unique stable fixed point there. With $m_Z(c)=(B+cI)^{-1}(b+c\mu_Z)$, after $r$ terminal calls from an entrance point,
\begin{equation}\label{eq:refine}
\norm{x_{h+r}-m_Z(c_L)}\le A_r:=6(3/4)^r+12e^{-1/(64\alpha)}.
\end{equation}
\end{lemma}
\begin{proof}
Squared-distance gaps and $\pi_j/\pi_Z<2$ give
\begin{equation}\label{eq:leak}
1-w_Z(x,v_L)\le4e^{-1/(33\alpha)},\qquad x\in Q_Z.
\end{equation}
The closest two templates have squared-distance gap at least $\alpha/16$. All remaining gaps exceed $1.89$. Each terminal coordinate update is a convex combination of the current coordinate and
\begin{equation}
\mu_{Z,j}+\frac{e_j+c_L(\bar\mu_j-\mu_{Z,j})}{B_{jj}+c_L}.
\end{equation}
The latter is within $.09$ of the template coordinate, proving invariance. Also $\norm{C}\le9(1-w_Z)$, so the terminal map is Lipschitz with constant less than $.751$ on each $Q_Z$, and has a unique stable fixed point there.

The exact recursion around $m_Z(c_L)$ has a linear part with norm at most $3/4$ and accumulated soft-responsibility error at most $12e^{-1/(33\alpha)}$. This yields~\eqref{eq:refine}. The conditional estimator $m_Z(c_L)$ is an analysis center, never a substitute for the actual iteration. Thus the bound controls the raw output without computing its terminal fixed point.
\end{proof}

\subsection{From conditional control to unconditional risk}
\begin{lemma}[A risk interface for the raw output]\label{lem:risk}
Under the continuous model of Theorem~\ref{thm:main}, any output of the observable controller with at least $r$ terminal calls on the good event obeys
\begin{equation}\label{eq:risk_interface}
R_{\rm coarse}\le R_{\rm oracle}+[\sqrt{A_r^2+T_\alpha}+\alpha^3]^2,
\end{equation}
where $T_\alpha=128\sqrt8e^{-1/(4\cdot10^6\alpha^2)}$. Restricted controllers obey the lower bound established below.
\end{lemma}
\begin{proof}
Let $E_\alpha=\{\norm{u}_\infty,\norm{\eta}_\infty\le\alpha/1000\}$. Then
\begin{equation}
\Prob(E_\alpha^c)\le8e^{-1/(2\cdot10^6\alpha^2)},
\end{equation}
and $E_\alpha$ implies~\eqref{eq:event}. For all observations, coordinatewise convexity of the update bounds the fourth moment uniformly in the call budget. Cauchy--Schwarz gives
\begin{equation}
\E[\norm{\widehat x-m_Z(c_L)}^2\mathbf1_{E_\alpha^c}]
\le T_\alpha:=128\sqrt8e^{-1/(4\cdot10^6\alpha^2)}.
\end{equation}
Conditioning on $(y,Z)$ and any independent controller random seed, the conditional mean is $m_Z(1)$ and orthogonality gives exactly
\begin{equation}
R(\widehat x)=R_{\rm oracle}+\E\norm{\widehat x-m_Z(1)}^2.
\end{equation}
The coefficient mismatch satisfies
\begin{equation}
m_Z(c_L)-m_Z(1)=(c_L-1)(B+c_LI)^{-1}[\mu_Z-m_Z(1)],
\end{equation}
whose $L^2$ norm is at most $\alpha^3$. Combining these bounds proves~\eqref{eq:risk_interface}.
For $r=r_\alpha$ and $\alpha\le10^{-5}$, $A_r\le\alpha^3$ and $2T_\alpha\le\alpha^6$, proving~\eqref{eq:risk}. On $Z=1,E_\alpha$, restricted iterates instead have error at least $.7-\alpha/1000$. Consequently
\begin{equation}
R_{\rm bad}\ge .30(.7-\alpha/1000)^2
 [1-8e^{-1/(2\cdot10^6\alpha^2)}]_+,
\end{equation}
which proves the limiting lower bound. Random switching does not change either argument: all estimates are pointwise before taking expectations.
\end{proof}
Lemmas~\ref{lem:barrier}--\ref{lem:risk} assemble Theorem~\ref{thm:main}: the first applies under arbitrary restricted switching, the second creates terminal entry using only observable decisions, and the last two convert entry into unconditional risk for the unaltered RED output.

\section{Preset Switching and Perturbation Stability}
\begin{proposition}[An observation-only preset]\label{prop:preset}
Compute $\widehat Z(b)$ once and set $h(0)=0$, $h(1)=25$, $h(2)=5$. Use any intermediate variances in $[.69,.71]$ for $h(\widehat Z)$ calls, then use $v_L$ up to $K=25+r$ calls. Under the event in~\eqref{eq:event}, switching occurs in $G_Z$. Under the Gaussian model and budget of Theorem~\ref{thm:main}, this preset attains the same $R_{\rm oracle}+5\alpha^6$ guarantee.
\end{proposition}
Theorem~\ref{geo:thm:main} improves this preset to $0/25/0$ using $\widetilde G_2$, including throughout the independent box. The present proposition retains the older entrance sets for comparison with the archived diagnostics.

The proof uses inward margins after first entry into $G_1$ and $G_2$. The preset uses those margins to finish its allocated calls. Entrance feedback can switch earlier and devote more of the budget to terminal refinement. Supplementary Section~\ref{sec:stability_proof} gives the face bounds for both protocols.

At the same $v=.70$, a shared switch can fail even when the preset works. Proposition~\ref{prop:shared_short} encloses a nonempty parameter box on which each common $h\in\{0,\dots,25\}$ leaves one class trapped on the wrong side under all subsequent terminal calls. The class that fails changes between $h=12$ and $h=13$.

\begin{proposition}[Uniform update perturbations]\label{prop:robust}
Replace $T_v(x_k)$ by $T_v(x_k)+d_k$ with a measurable perturbation based on observations and history. Assume $\norm{d_k}_\infty\le10^{-4}$ globally, and $\norm{d_k}_\infty\le\delta_L\le10^{-4}$ on terminal calls. Keep the observation rule, initialization and scale choices. On~\eqref{eq:event}, the restricted class-1 trajectories still satisfy $s_k\le.3$, and both intermediate-band protocols retain their $0/25/5$ entry bounds. On the Gaussian model, with $K=25+r$, their raw outputs satisfy
\begin{align}
\widetilde R-R_{\rm oracle}
&\le\big[\sqrt{(A_r+4\sqrt{2}\delta_L)^2+T_\alpha}
             +\alpha^3\big]^2. \label{eq:robust_risk}
\end{align}
\end{proposition}
For a denoiser approximation the update error is $\gamma\lambda(\widetilde D_\sigma-D_\sigma)$, so these bounds apply to the corresponding scaled coordinatewise error. With $r=r_\alpha$, a terminal error $o(\alpha^2)$ preserves the asymptotic oracle ratio, and $O(\alpha^3)$ preserves the $O(\alpha^6)$ excess-risk order. One concrete application allows each denoiser log weight to move by at most $10^{-4}$ at fixed true distribution, templates and variances. The intermediate update error then meets the constant radius, and terminal posterior concentration makes its error exponentially small on the entrance regions. The global bound still controls the exceptional observations. See Supplementary Section~\ref{sec:stability_proof} for the bounds.

\paragraph{Terminal coefficient calibration.}
For a positive coefficient $c$, the conditional analysis center satisfies
\begin{align}
&\E\norm{m_Z(c)-m_Z(1)}^2 \nonumber\\
&\quad=\alpha^4(c-1)^2
\left[\frac{\alpha}{(\alpha+c)^2(1+\alpha)}
      +\frac{1}{2(1+c)^2}\right].
\label{eq:coefficient_exact}
\end{align}
Indeed, in a coordinate with eigenvalue $\beta$ of $B$, the variance of $(b-B\mu_Z)_j$ is $\alpha^4\beta(1+\beta)$. Subtracting the two linear centers and summing over $\beta=\alpha,1$ gives~\eqref{eq:coefficient_exact}. For fixed $c>0$, the center has risk ratio tending to $1+(c-1)^2/[3(c+1)^2]$. Thus the order $\alpha^4$ within a correctly selected component and the oracle ratio involve different calibration requirements. The formula is for the center. Proposition~\ref{prop:robust} uses the original coefficient window to control the actual iterates.

\section{A Settled Common Platform as a Contrast}\label{sec:platform}
The main controller switches when a transient trajectory reaches an observable entrance. A different protocol holds a single variance $v$ common to all observations long enough to approach equilibrium. Its local scale requirement explains why the stopping rule matters.

Assume $|v-v_0|\le10^{-3}$ and hold the platform for at least
\begin{equation}
N_{\rm settle}=\left\lceil\frac{\log(20/\alpha)}{\log(25/23)}\right\rceil.
\end{equation}
\begin{proposition}[Local necessity for a settled common platform]\label{prop:calibration}
On~\eqref{eq:event}, if $v-v_0\ge20\alpha$, class 1 enters a terminal incorrect invariant region. If $v_0-v\ge20\alpha$, class 0 does so. In either case template error remains greater than $.5$ after any number of terminal steps. Therefore simultaneous recovery requires $|v-v_0|<20\alpha$ within this protocol.
\end{proposition}
The failure premise is nonempty only when $\alpha\le5\cdot10^{-5}$. It does not apply to short transient platforms or state-dependent switching.

\paragraph{Mechanism of the necessity result.}
On $\mathcal R=[-.1,.75]\times[-.1,.15]$, $H_v=Dg_v$ obeys $.32I\preceq H_v\preceq2.01I$ and $\norm{DT_v}\le.92$. The virtual fixed point $r_*(v)$ has $(r_*(v_0))_1=.5$. The derivative signs $\partial_vg_1>.12$, $\partial_vg_2<-.16$, and $H_{12}>0$ imply
\begin{equation}
-\frac{d(r_*)_1}{dv}
=\frac{H_{22}\partial_vg_1-H_{12}\partial_vg_2}{\det H}>.1.
\end{equation}
Actual class-0/1 fixed points differ from the virtual point by at most $1.572\alpha$, and settlement leaves at most $\alpha/10$ error. A $20\alpha$ offset consequently places one class more than $.32\alpha$ across the incorrect side of the terminal separator, inside an incorrect invariant region. Supplementary Section~\ref{sec:platform_proof} supplies the terminal trapping argument; the conclusion is actual permanent failure, not a reversed sufficient condition.

\paragraph{A separate sufficient direction.}
Proposition~\ref{prop:common} states that $|v_k-v_0|\le\alpha/1000$ for
\begin{equation}
H_+(\alpha)=\left\lceil\frac{\log(64/\alpha)}{\log(25/23)}\right\rceil
\end{equation}
intermediate calls suffices for correct terminal entry. Its prescribed-window crossing lower and upper bounds are in Proposition~\ref{prop:cross}. This supplies the successful direction alongside the local necessary condition, with different constants and explicitly different holding-time requirements. It does not make every scale inside the necessary window sufficient. Combining $H_+$ with $r_\alpha$ terminal calls gives 317 calls at $\alpha=10^{-5}$, while the observable controller needs the sufficient budget $25+r_\alpha=154$. Both use the same terminal objective. The logarithmic crossing lower bound belongs to the shared-window protocol, not to all controllers in Table~\ref{tab:protocols}.


\section{Cross-geometry and photo-patch diagnostics}
\begin{figure*}[t]
\centering
\includegraphics[width=\textwidth]{../../experiments/results/geometry_prediction_wide/geometry_prediction.pdf}
\caption{Four tests on identical geometry coordinates. Dark means the
test holds. The first two panels use 1200 logarithmically sampled
variances between $v_L$ and 100, plus the analytic infinite-variance
endpoint for the first panel. They concern the fixed class-1 corner
$(.3,-.01)$ and noiseless data. The third panel checks the analytic
class-1 channel conditions for $[.69,.71]$ and 25 calls, with
$\delta=10^{-4}$. Light means unclassified by those sufficient
conditions. The fourth panel checks exact 106-call iterates for all
three noiseless template observations, with radius-$.1$ success.
All panels fix $\alpha=.001,a=1,r=1.29$. The grid is logarithmic in
$h$ and linear in $d$. It was widened after an initial narrower grid
showed no negative points in the first two tests. The grid defines
a diagnostic sample and supplies no continuous-parameter guarantee.}
\end{figure*}
The four masks have respectively 995, 1673, 1383 and 1248 positive
points out of 1681. Every channel-certified point enters the class-1
prepared rectangle. Only 1176 of these points recover all three
classes. Other classes' entry, terminal distance gaps and available
refinement budget require separate tests as specified in
Theorem~\ref{thm:criterion}.

\subsection{Matched risk illustrations}
\begin{table}[t]
\centering\small
\setlength{\tabcolsep}{3pt}
\begin{tabular}{lrrrr}
\toprule
Model & $K$ & Preset & Restricted & Terminal\\
\midrule
Calibrated & 106 & $1.53\cdot10^{-12}$ & $.3163$ & $.3163$\\
Separate noise & 58 & $1.28\cdot10^{-8}$ & $.3169$ & $.3169$\\
Fine calibration & 154 & $1.47\cdot10^{-20}$ & $.3169$ & $.3169$\\
\bottomrule
\end{tabular}
\caption{Prior-weighted squared error of the raw RED iterates on paired draws.
All methods within a row share the observation, initialization, total
budget and terminal variance.}
\label{tab:risk_diagnostics}
\end{table}
We fix $a=1$, $d=2$, $r=1.29$ and $h=1.78$ in all three rows of
Table~\ref{tab:risk_diagnostics}. The calibrated row uses
$\alpha=10^{-3}$, $\tau=\nu=\alpha^2$ and $\theta=0$.
The second row uses $\alpha=10^{-6}$, $\tau=10^{-4}$,
$\nu=10^{-5}$, $v_L=(\tau^2+\nu^2)^{1/2}$ and $\lambda=1$.
The preset uses $v=.70$ for its observation-selected intermediate calls.
The restricted policy alternates $.17$ and $2.5$ for 25 calls, then
uses $v_L$. The terminal comparator uses $v_L$ throughout.
The fine row repeats the calibration at $\alpha=10^{-5}$.
Budgets follow the respective refinement formulas.

We draw 512 independent pairs per class per row using seed 20261109
and weight class averages by $\pi$. All settings are prespecified.
All draws satisfy the coordinate bounds in~\eqref{eq:ind_event}.
In all three rows, mode accuracy and radius-$.1$ template-neighborhood success
are both $1$ for the preset and $.683$ for the two comparators.
We compute these metrics separately from squared error.
The table illustrates the separation.
The guarantees for the continuum of variances and arbitrary restricted
switching follow from the common barrier proof.


\paragraph{Fine excess risk.}
For the fine row, the analytic oracle risk is $1.49999\cdot10^{-20}$.
We estimate the excess risk directly as the prior-weighted mean of
$\norm{\widehat x-m_Z(1)}^2$, which is $8.49\cdot10^{-32}$ in double precision. This is an implementation
reading near numerical resolution, dominated by rounding. Its specific
value does not measure agreement with the $5\cdot10^{-30}$ bound.
Repeating the intermediate steps at 50 and 80 digits gives
$3.48001\cdot10^{-34}$ at both precisions using exact decimal parameters
and the same binary64 normal draws. After verifying terminal entry for
every sample, we propagate the affine residual
$[I-(B+I)/4]^n(x_h-m_Z(1))$.
The omitted terminal leakage changes the output by at most
$12e^{-.03/v_L}<10^{-130288343}$.
Here $m_Z(1)=(B+I)^{-1}(b+\mu_Z)$; labels enter evaluation only. The ordinary empirical MSE fluctuates on a
larger scale than this excess.

\paragraph{Precision reference.}
The program \path{experiments/code/precision_diagnostic.py} replays
the saved normal draws at 50 and 80 decimal digits. Parameters are
exact decimal constants; the standard-normal draws are the same
binary64 numbers in the original sampling stream. For each of 1536
samples it checks terminal entry before using the analytic affine
recursion. The terminal error bound $12\exp(-.03/v_L)$ comes from
the uniform distance gap $.06$ and rectangle invariance. Its
base-10 logarithm is less than $-130288343$ here. Both precisions
give excess risk $3.48001105532541\cdot10^{-34}$.
This controls the omitted soft-template term but does not constitute
an interval enclosure of the arbitrary-precision arithmetic.

\paragraph{Initialization and weak-coefficient boundaries.}
The least-squares start lies inside the correct terminal rectangle
when $|e_1|\le .1\alpha$ and $|e_2|\le.005$.
For exact updates with $\lambda=\alpha/4$, class-1 drift at $s=.3$
is strictly outward on $E$ for every admissible scale, since
$b_1-\alpha s\ge .45\alpha$ and
$c(v)(m_v-s)\ge-(1.02+.3)\alpha/4$.
Changing the data loss to a likelihood-normalized form would also
change the drift and step-size conditions.

\paragraph{Weak-regularization budget.}
For noiseless observations of the three nominal templates,
$\alpha=10^{-5}$ and $\lambda=\alpha/4$, terminal-only iteration
from the backprojection has prior-weighted squared error $.75362$,
$.74998$, $.71238$, $.42955$ and $.001932$ at budgets
154, $10^3$, $10^4$, $10^5$ and $10^6$.
The first radius-$.1$ entry times are 1, 828881 and 736823.
These are observed trajectories at a changed coefficient, not
uniform hitting-time guarantees. See
\path{experiments/code/weak_budget_diagnostic.py}.

\paragraph{Patch data and protocol.}
The source images are the two photos distributed in
\texttt{sklearn.datasets.load\_sample\_images}.
They are converted to grayscale and divided into nonoverlapping
$4\times4$ patches. A seed-20261110 subset of 2048 patches from
\texttt{china.jpg} trains eight means with $k$-means (five starts);
the masses are cluster frequencies and the common isotropic variance
is the mean squared training residual per coordinate.
A disjoint-photo subset of 256 patches from \texttt{flower.jpg}
is evaluated. No labels are defined for these empirical truths.
The exact MMSE denoiser of the fitted mixture uses total variance
$\tau^2+\sigma^2$ and the RED residual is unnormalized.
The blur applies periodic convolution by $[.25,.5,.25]$ on each
patch axis. This is a patch-level degradation operator.
All policies share the resulting noisy observations, backprojection,
154 calls, $\lambda=.1$, $\gamma=.25$ and terminal $\sigma=.02$.
Their mean squared Euclidean errors are respectively $.00990445$,
$.00988671$ and $.00986463$ for terminal, intermediate, and low/high
protocols. Per-pixel MSE divides these by 16. These are descriptive
results on one image pair with no schedule search or fitted-model
certificate. See \path{experiments/code/patch_diagnostic.py}.

Reproduction uses the repository's Python environment and a fresh
output directory for each of
\path{experiments/code/geometry_prediction_diagnostic.py},
\path{experiments/code/precision_diagnostic.py}, and
\path{experiments/code/patch_diagnostic.py}, each with
\texttt{--out <new-directory>}. Dependencies are NumPy, SciPy,
Matplotlib, Pillow, scikit-learn and mpmath.

\section{Operator-aware terminal contraction for finite mixtures}
\label{sec:terminal_history}
This section retains the original RED update and permits any finite
number of equal-width isotropic Gaussian components, any ambient
dimension, and any linear observation operator.
Write $m_v=\sum_iw_i\mu_i$ and
$C_v=\sum_iw_i(\mu_i-m_v)(\mu_i-m_v)^\top$.
This is the posterior covariance of the template means.
It excludes the within-component covariance.
Throughout, $v>\tau^2$, $c=\lambda(1-\tau^2/v)>0$ and $B=A^\top A$.
Conditional mean--covariance derivative identities have an established
foundation~\cite{dytso2023}. Our use is a model-input certificate for
the original RED terminal map.

\subsection{An exact local test using the sensing operator}
Let $\mathcal S=\operatorname{span}\{\mu_i-\mu_0\}$ and let
$U,V$ be orthonormal bases of $\mathcal S$ and its orthogonal
complement. Define
\begin{align}
 G_c={}&U^\top BU+cI \nonumber\\
 &-U^\top BV(V^\top BV+cI)^{-1}V^\top BU,\\
 C_S&=U^\top C_vU,\qquad
 \rho_v(x)=\frac cv\lambda_{\max}(G_c^{-1/2}C_SG_c^{-1/2}).
\end{align}
The term involving $V$ vanishes when $\mathcal S$ is the whole space.
For coincident means set $\rho_v=0$.
\begin{proposition}[Operator--posterior curvature test]\label{prop:operator_posterior}
The Hessian of the exact RED potential is
\[
 H_v=B+cI-(c/v)C_v.
\]
At each state, $H_v\succ0$ if and only if $\rho_v(x)<1$.
On a convex set where $\rho_v\le\bar\rho<1$,
\[
 (1-\bar\rho)(\lambda_{\min}(B)+c)I
 \preceq H_v\preceq(\lambda_{\max}(B)+c)I.
\]
For $0<\gamma\le(\lambda_{\max}(B)+c)^{-1}$,
the map is Lipschitz with factor
$q=1-\gamma(1-\bar\rho)(\lambda_{\min}(B)+c)<1$ on that set.
\end{proposition}
\begin{proof}
Differentiating the responsibilities gives $D_xm_v=C_v/v$.
In the basis $[U,V]$, the lower-right block of $H_v$ is positive
definite and its Schur complement is $G_c-(c/v)C_S$.
Moreover $U^\top(B+cI)^{-1}U=G_c^{-1}$, so the nonzero eigenvalues
in the definition of $\rho_v$ equal those of
$(c/v)(B+cI)^{-1/2}C_v(B+cI)^{-1/2}$.
The Loewner bounds follow.
Integrating $I-\gamma H_v$ on a line segment in the convex set
gives the stated Lipschitz factor.
\end{proof}
Constructing $G_c$ requires solves involving the actual operator.
Only the final covariance eigenproblem has dimension at most $K-1$.
Pairwise secant gains alone do not lower-bound $B$ on $\mathcal S$.
For example, means $(0,0),(1,0),(2,1)$ under
$B=\operatorname{diag}(1,0)$ have pair gains $1,4/5,1/2$ but
span a null direction.

\subsection{A common invariant terminal ball}
Fix the terminal map and a ball
$\mathcal B=\{x:\norm{x-\bar x}\le R\}$.
For reference component $j$, set $d_{ij}=\mu_i-\mu_j$ and
\begin{align}
 \Delta_{ij}&=\norm{\bar x-\mu_i}^2-\norm{\bar x-\mu_j}^2,\\
 \beta_{ij}&=\operatorname{logistic}\left(
 \log\frac{\pi_i}{\pi_j}-\frac{\Delta_{ij}-2R\norm{d_{ij}}}{2v_L}\right),\\
 \overline C_{\mathcal B}&=\sum_{i\ne j}\beta_{ij}d_{ij}d_{ij}^\top.
\end{align}
The affine distance-gap bound implies $w_i/w_j\le
\exp(\operatorname{logit}\beta_{ij})$, hence $w_i\le\beta_{ij}$.
Since
$C_v=\sum_iw_id_{ij}d_{ij}^\top-(m_v-\mu_j)(m_v-\mu_j)^\top$,
this is a valid covariance upper bound on the whole ball.

\begin{proposition}[Common terminal ball]\label{prop:common_terminal_ball}
Let $\overline C$ be any whole-ball covariance upper bound and put
$M=B+c_LI$. Suppose $(c_L/v_L)\overline C\preceq\bar\rho M$
for some $\bar\rho<1$, the step condition of
Proposition~\ref{prop:operator_posterior} holds, and
\[
 \norm{T_{v_L}(\bar x)-\bar x}\le(1-q)R.
\]
Then the ball is invariant and contains a unique terminal fixed point.
For any two starts in the ball at common time $H$,
\[
 \norm{x^{(a)}_{H+n}-x^{(b)}_{H+n}}
 \le q^n\norm{x^{(a)}_H-x^{(b)}_H},\qquad n\ge0.
\]
\end{proposition}
\begin{proof}
Contraction gives
$\norm{T(x)-\bar x}\le q\norm{x-\bar x}+\norm{T(\bar x)-\bar x}\le R$.
The contraction mapping theorem applies on the closed ball.
Iteration of the Lipschitz estimate proves the distance bound.
\end{proof}
If both updates have norm errors bounded by $\delta_a,\delta_b$
and remain in the ball, the distance bound gains
$(\delta_a+\delta_b)(1-q^n)/(1-q)$.
For perturbed invariance add the corresponding error bound to the
center residual.
An almost-sure certificate with a uniform $q$ implies
\[
 |\sqrt{R_a(n)}-\sqrt{R_b(n)}|
 \le q^n\big(\mathbb E\norm{x_H^{(a)}-x_H^{(b)}}^2\big)^{1/2}.
\]
A good-event-only certificate also needs an exceptional-event
$L^2$ term. A finite batch does not imply a population risk guarantee.
Neither common convergence nor a unique ball fixed point implies
accurate recovery.

\subsection{Sharper pairwise and global covariance bounds}
For each pair let
\begin{align}
 t_{ij}&=\log(\pi_i/\pi_j)
 -\frac{\norm{\bar x-\mu_i}^2-\norm{\bar x-\mu_j}^2}{2v},\\
 L_{ij}&=t_{ij}-R\norm{\mu_i-\mu_j}/v,\nonumber\\
 U_{ij}&=t_{ij}+R\norm{\mu_i-\mu_j}/v .
\end{align}
Let $\alpha_i=(\sum_k e^{L_{ki}})^{-1}$ and let $t^\circ_{ij}$
be the projection of zero onto $[L_{ij},U_{ij}]$.
Every state in the ball has $w_i\le\alpha_i$ and
\[
 w_iw_j\le\kappa_{ij}
 =\min\left\{\alpha_i\alpha_j,
 \frac1{(1+e^{t^\circ_{ij}})(1+e^{-t^\circ_{ij}})}\right\}.
\]
The second bound follows by maximizing the product at fixed pair
odds and total mass at most one.
The exact identity
$C_v=\sum_{i<j}w_iw_j(\mu_i-\mu_j)(\mu_i-\mu_j)^\top$
then gives a whole-ball bound by replacing $w_iw_j$ with $\kappa_{ij}$.
One may choose the valid covariance envelope with smaller normalized
spectral bound. An elementwise minimum of matrices is unjustified.

A separate global sufficient bound uses the whitened diameter
$D_M^2=\max_{i,j}(\mu_i-\mu_j)^\top M^{-1}(\mu_i-\mu_j)$.
The directional bounded-variable variance inequality yields
$\rho_v(x)\le cD_M^2/(4v)$ everywhere.
If this bound is below one, the terminal map is globally contractive
at the prescribed step size.
Failing a sufficient bound leaves the global question open.

\subsection{Certificate for the fixed patch batch}
The template span in the photo-patch case has dimension seven.
Its restricted sensing spectrum has minimum $.02509617$,
compared with the pairwise secant minimum $.22752$.
The whitened-diameter global bound is $1.21022$, so it is inconclusive.
For each observation we compute a candidate center by ten Newton
steps from its 5000-call terminal state.
At common time 154 the radius is 1.01 times the maximum distance
of the three saved outputs from this center, with floor $10^{-9}$.
No truth patch is used in this construction.
The original reference-component envelope passes 255 of 256 balls.
The pairwise envelope passes all 256.

The program \path{experiments/code/patch_terminal_interval.py}
uses 60-digit outward-rounded interval arithmetic to check all 256.
It verifies the Gram decomposition with exact rational arithmetic,
proves positive definiteness of
$.45(B+c_LI)-(c_L/v_L)\overline C_{\mathcal B}$ by interval LDL,
and checks center residuals, coverage and step size.
Since $B\succeq0$ and $\norm{B}=1$, it yields the uniform factor
$q=1-\gamma(.55c_L)<.999429$.
The certificate treats exported binary64 coefficients and snapshots
as exact inputs to the subsequent exact terminal map.
It does not enclose the earlier floating-point propagation or
subsequent finite-precision perturbations.
It proves neither a new-observation risk bound nor low error of the
shared fixed points.

\subsection{Related contraction and recovery results}
Sinha et al.~\cite{sinha2025} study fixed linear kernel denoisers,
including appropriately scaled PnP for nonsymmetric kernels.
The present covariance varies with state and the update is residual-additive RED.
Wolf et al.~\cite{wolf2026} jointly analyze denoising covariance
and generalized backprojection in PnP forward--backward splitting.
Their pointwise and distributional bounds use support and restricted
conditioning assumptions. Their compact-support assumptions and
update do not directly apply to our untruncated GMM RED model.
The covariance identities, Schur complements and contraction tools
are established mathematics. The present statements organize them
as model-input checks and a batch-specific terminal certificate.

\section{Coefficient and posterior changes in a scale update}
\label{sec:coefficient_split}
Write $T_{c,v}(x)=x-\gamma[Bx-b+c(x-m_v(x))]$.
The original algorithm remains on $c=\lambda(1-\tau^2/v)$.
For fixed $x$ the exact identity is
\begin{align}
 T_{c,v}(x)-T_{\tilde c,\tilde v}(x)
 ={}&-\gamma(c-\tilde c)(x-m_v(x))\nonumber\\
 &+\gamma\tilde c(m_v(x)-m_{\tilde v}(x)).
\end{align}
For $z=1/v$ and $E_i(x)=\norm{x-\mu_i}^2/2$,
\begin{align*}
 \partial_z m_{1/z}(x)&=-\operatorname{Cov}_{w}(\mu,E),\\
 \norm{\operatorname{Cov}_w(\mu,E)}
 &\le\sqrt{\lambda_{\max}(C_w)\operatorname{Var}_w(E)} .
\end{align*}
Both formulas follow by differentiating the normalized exponential
weights and applying directional Cauchy--Schwarz.
Integration bounds $\norm{m_v-m_{\tilde v}}$ by
$|v^{-1}-\tilde v^{-1}|$ times the supremum of this covariance norm.
On a region containing both trajectories, if the baseline map is
$q_k$-Lipschitz, let
$R_k=\sup_x\norm{x-m_{v_k}(x)}$ and
$M_k=\sup_{x,z}\norm{\operatorname{Cov}_{w(x,1/z)}(\mu,E(x))}$.
Then their distance satisfies
\begin{align}
 A_\ell&=|c_\ell-\tilde c_\ell|R_\ell+
 \tilde c_\ell|v_\ell^{-1}-\tilde v_\ell^{-1}|M_\ell,\nonumber\\
 e_n&\le\left(\prod_{k=0}^{n-1}q_k\right)e_0
 +\gamma\sum_{\ell=0}^{n-1}
 \left(\prod_{k=\ell+1}^{n-1}q_k\right)A_\ell .
\end{align}
This is the iterated one-step comparison.
It permits cancellation of vector effects but does not identify
unique statistical contributions from a single measured output difference.

There is an explicit interaction when two separate ablations are
both measured from the corner $(c_0,v_0)$.
For $m_i=m_{v_i}(x)$,
\begin{align}
 &T_{c_1,v_1}-T_{c_1,v_0}-T_{c_0,v_1}+T_{c_0,v_0}\nonumber\\
 &\hspace{18mm}=\gamma(c_1-c_0)(m_1-m_0).
\end{align}
Thus coefficient-only and posterior-only changes do not generally
sum to the combined change. Its interaction norm is at most
$\gamma|c_1-c_0|\norm{m_1-m_0}$, before state propagation.

\section{Several competing templates and a mass threshold}
\label{sec:multi_competitor}
This result extends the fixed-state three-template comparison.
All components have the common isotropic covariance $\tau^2I$.
Let $\mu_0=0$, $\mu_1=e$ with $\norm e=1$, and fix $x$ with
$s=e^\top x\in(0,1/2)$ and $P=1/2-s$.
For $M\ge1$ competitors $\nu_j$, suppose
\begin{align}
 e^\top\nu_j&=-a_j,\qquad a_j>0,\nonumber\\
 Q_j&=\tfrac12(\norm{x-\nu_j}^2-\norm x^2)>P.
 \label{eq:multi_geometry}
\end{align}
Their prior ratios are $1:r:hq_1:\cdots:hq_M$ where
$r,h,q_j>0$ and $\sum_jq_j=1$.
The internal composition $q$ remains fixed as the total ratio $h$
varies. Define
\begin{align}
 \bar a&=\sum_jq_ja_j,\quad
 \bar Q=\sum_jq_jQ_j,\quad
 \overline{aQ}=\sum_jq_ja_jQ_j,\nonumber\\
 C_s&=(1-s)r-s>0,\nonumber\\
 h_{\max}&=\frac{C_s}{\bar a+s},\qquad h_0=\frac r{\bar a}.
 \label{eq:multi_mass_range}
\end{align}
Let $z_L=1/v_L$ with
\begin{equation}
 v_L>\tau^2,\qquad r e^{-Pz_L}\le \frac{s}{1-s}.
 \label{eq:multi_endpoint_conditions}
\end{equation}
For inverse variance $z$, put $l(z)=1-\tau^2z$ and
\begin{align}
 m_h(z)&=
 \frac{r e^{-Pz}-h\sum_j a_jq_je^{-Q_jz}}
      {1+r e^{-Pz}+h\sum_jq_je^{-Q_jz}},\nonumber\\
 g_h(z)&=l(z)(m_h(z)-s).
 \label{eq:multi_drift}
\end{align}
The projected RED displacement per unit step size is
$e^\top(b-Bx)+\lambda g_h(z)$.
Thus comparison of $g_h$ compares actual updates at this state,
for the same $\lambda>0$, without alignment of $e$ and $B$.

\begin{proposition}[Mass threshold with several competitors]
\label{prop:multi_mass}
Under \eqref{eq:multi_geometry}--\eqref{eq:multi_endpoint_conditions},
for $0<h\le h_{\max}$, the infinite-variance endpoint
$g_h(0)$ dominates the terminal endpoint $g_h(z_L)$.
A finite intermediate variance strictly improves on both endpoints
if and only if $D(h)>0$, where
\begin{align}
 D(h)={}&-Pr\nonumber\\
 &+h[(1+r)\overline{aQ}+r\bar Q-Pr(1+\bar a)]\nonumber\\
 &+h^2(\overline{aQ}-\bar a\bar Q)\nonumber\\
 &-\tau^2[C_s-h(\bar a+s)](1+r+h).
 \label{eq:multi_polynomial}
\end{align}
There is exactly one root $h_{\rm crit}\in(0,h_0)$.
On $0<h\le h_{\max}$, improvement is equivalent to
$h>h_{\rm crit}$. At fixed remaining scalar inputs, increasing
any $Q_j$ strictly decreases $h_{\rm crit}$.
\end{proposition}

\begin{proof}
The mass range gives
\[
 u:=g_h(0)=\frac{C_s-h(\bar a+s)}{1+r+h}\ge0.
\]
Without the competitors, the terminal right-template responsibility
is at most $s$ by \eqref{eq:multi_endpoint_conditions}.
Adding negative-projection templates decreases its horizontal mean,
so $g_h(z_L)\le0$ and the high-variance endpoint is larger.

Set
\begin{align*}
 A_u(z)&=(1-s)l(z)-u,\\
 U_u(z)&=\frac{s l(z)+u}{A_u(z)},\\
 W_{u,j}(z)&=\frac{(a_j+s)l(z)+u}{A_u(z)}.
\end{align*}
Since $u\le C_s/(1+r)<1-s$, $A_u(0)>0$.
The set $I_u=\{z\in[0,z_L]:A_u(z)>0\}$ is an initial interval.
On that interval, multiplication by the positive responsibility
denominator gives the equivalence
\begin{align}
 g_h(z)>u
 &\quad\Longleftrightarrow\quad J_u(z)<r,\nonumber\\
 J_u(z)&=U_u(z)e^{Pz}
     +h\sum_jq_jW_{u,j}(z)e^{-(Q_j-P)z}.
 \label{eq:multi_convex_test}
\end{align}
In particular $J_u(0)=r$.
The identities
\begin{align*}
 U_u&=\frac{s}{1-s}+\frac{u}{(1-s)A_u},\\
 W_{u,j}&=\frac{a_j+s}{1-s}
             +\frac{(1+a_j)u}{(1-s)A_u}
\end{align*}
show that each ratio is log-convex.
Indeed a positive affine reciprocal is log-convex on its positive
domain, and positive sums preserve log-convexity.
Multiplication by an exponential preserves that property.
Moreover $U_u>0$, $U_u',U_u''\ge0$ and $P>0$, so the first
term in $J_u$ is strictly convex. Thus $J_u$ is strictly convex.

No improvement is possible where $A_u(z)\le0$.
All template projections are at most 1 and the weight on $\mu_0$
is positive, giving $m_h(z)<1$. Since $l(z)>0$,
\[
 g_h(z)-u<(1-s)l(z)-u=A_u(z)\le0.
\]
Consequently an improvement anywhere in $(0,z_L)$ exists exactly
when $J_u'(0)<0$. This argument does not require $A_u$ to remain
positive all the way to $z_L$.

On $I_u$ the exact identity
\[
 g_h(z)-u=
 \frac{A_u(z)e^{-Pz}[r-J_u(z)]}
      {1+r e^{-Pz}+h\sum_jq_je^{-Q_jz}}
\]
shows that the derivatives of $g_h$ and $J_u$ at zero have
opposite signs. Direct differentiation of
\eqref{eq:multi_drift} yields
\[
 g_h'(0)=\frac{D(h)}{(1+r+h)^2},
\]
which proves the improvement criterion.

For uniqueness,
$D(0)=-Pr-\tau^2C_s(1+r)<0$.
At $h_0=r/\bar a$, $m_{h_0}(0)=0$ and
\[
 m_{h_0}'(0)=
 \frac{r(\overline{aQ}/\bar a-P)}{1+r+h_0}>0.
\]
Hence $g_{h_0}'(0)=m_{h_0}'(0)+\tau^2s>0$ and $D(h_0)>0$.
A polynomial of degree at most two with opposite signs at these
endpoints has exactly one root inside, which is simple and crosses
from negative to positive. Since $h_{\max}<h_0$, this proves the
claimed threshold on the comparison range.
Finally,
\[
 \frac{\partial D(h)}{\partial Q_j}
 =h q_j[a_j(1+r+h)+(r-h\bar a)]>0
 \quad(0<h<h_0).
\]
At the root $\partial_hD>0$, so implicit differentiation gives
$\partial_{Q_j}h_{\rm crit}<0$.
\end{proof}

The comparison of a partial derivative in $Q_j$ keeps $a_j$
and the other scalar inputs fixed. A geometric perturbation can
change several of these inputs at once.
For $M=1$, the covariance term in \eqref{eq:multi_polynomial}
vanishes and the three-template polynomial is recovered.
For $M>1$, $\operatorname{Cov}_q(a,Q)$ can be negative, so
global uniqueness on the positive $h$ axis is not asserted.

\paragraph{A realizable second positive root.}
Take $x=(.3,3)$, $e=(1,0)$, $r=1.29$,
$\tau^2=10^{-4}$, $v_L=.01$ and $q_1=q_2=1/2$, with
\begin{align*}
 \nu_1&=(-.1,3+\sqrt{28.93}),\\
 \nu_2&=(-3,3+\sqrt{2.2}).
\end{align*}
Then $a=(.1,3)$ and $Q=(10,2)$ give
\[
 D(h)=-5.799815h^2+15.09746335h-.258138087.
\]
Its positive roots are approximately $.01721192$ and $2.58588204$,
while $h_{\max}\simeq.32594595<h_0\simeq.83225806$.
Only the first root belongs to the threshold interval.

\paragraph{Scope and numerical checks.}
The result is a fixed-state single-step comparison.
It does not provide a channel, a terminal entrance, or a raw-output
risk separation for arbitrary competitor configurations.
Strict improvement over the infinite-variance endpoint persists
against a sufficiently large finite high variance by continuity.
The program \path{experiments/code/theory_extension_checks.py}
checks the derivative identity and convex criterion on 128 finite
cases, including 69 cases outside the former sufficient condition
$C_s/(1+r)<(1-s)l(z_L)$.
It also checks the explicit two-root geometry, the single-competitor
reduction and 64 four-corner attribution identities.
These checks are distinct from the external draft's reported
1500-case batch and do not replace the analytic proof.

\section{Frozen independent-noise and held-out-photo validation}
\label{sec:independent_patch}
This validation follows the original diagnostic and theory
development. Before evaluating its new observations, we export the
original china-trained eight-component prior and blur, and freeze
the restoration and ball-construction rules.
The saved \path{experiments/results/patch_independent_validation/plan.json}
records the design.
We reuse the original 256 flower truth patches with new noise.
We also take 256 nonoverlapping grayscale $4\times4$ patches from
each of Set12 images 01, 02 and 10, obtained from the KAIR testset.
These are the cameraman, house and boat photographs.
Image indices were chosen before restoration outcomes were available.
Patch selection uses seed $20261115+i$ for photo index $i=0,1,2,3$,
except that the original flower patches are retained.
Each photo receives three independent Gaussian noise realizations
with seed $20261112+j+1000i$, $j=0,1,2$.
Noise standard deviation remains $.01$.
There is no model refitting or schedule selection.

All original protocols retain $\lambda=.1$, $\gamma=.25$,
terminal $\sigma=.02$ and backprojection initialization.
The intermediate preset uses $\sigma=.1$ for 25 calls, and the
low/high comparator alternates $\sigma=.02/.3$ for 25 calls.
Each then uses terminal refinement.
We also repeat the coefficient-only and smoothing-only ablations.
These ablations intentionally alter the update for attribution.
Reported budgets are 25, 154, 1000 and 5000 calls.
Within each observation, all methods share the truth, noise,
initialization and total budget.

\begin{table}[t]
\centering\small
\setlength{\tabcolsep}{3pt}
\begin{tabular}{lrrrr}
\toprule
Photo & Terminal & Intermediate & Low/high & $\Delta$ (\%)\\
\midrule
Flower & .0098413 & .0098278 & .0098044 & $-.138$\\
Cameraman & .0235398 & .0235585 & .0235378 & $+.079$\\
House & .0135999 & .0136992 & .0136210 & $+.730$\\
Boat & .0214076 & .0213269 & .0213542 & $-.377$\\
\bottomrule
\end{tabular}
\caption{Mean squared Euclidean error at 154 calls, averaged over
256 truth patches and three noise realizations per photo.
$\Delta$ is the intermediate error change relative to terminal-only.
Negative means lower error. There are no true component labels for
the photo patches.}
\label{tab:independent_photo_risk}
\end{table}

The intermediate direction is consistent over the three noise
seeds within each image but differs across images.
The three new photos' equal-weighted mean errors are
$.01951577$, $.01952817$ and $.01950431$ for terminal,
intermediate and low/high, respectively.
These descriptive results show no consistent intermediate benefit.
Noise repeats share their truth patches, so the 3072 observations
are not 3072 independent photographs.
We report each photo and its seed-level differences separately,
without a population significance claim.
Later iteration generally increases error after 154 calls even
as the outputs of different schedules approach one another.

\begin{figure*}[t]
\centering
\includegraphics[width=\textwidth]{../../experiments/results/patch_independent_validation/risk_curves.pdf}
\caption{Frozen-model validation. Each panel averages three noise
realizations of the same 256 patches from one photo. Recovery error
and terminal convergence need separate evaluation. Smoothing-only
retains the terminal coefficient and varies posterior weights.
Coefficient-only retains terminal posterior weights and varies the
coefficient. These two curves are attribution controls.}
\label{fig:independent_photo_risk}
\vspace{6pt}
\includegraphics[width=\textwidth]{../../experiments/results/patch_independent_validation/distance_curves.pdf}
\caption{Intermediate versus terminal output distance and the
numerically evaluated sample-specific exact-map bound
$N^{-1}\sum_iq_i^{2(n-154)}d_{i,154}^2$.
The bound constants are interval verified.
The observed curves use floating-point iterates, whose per-step
rounding is not enclosed. They are empirical consistency checks,
and the centers use later states, so this is not an a priori
prediction of the 5000-call outputs.}
\label{fig:independent_distance}
\end{figure*}

For each observation, a candidate ball uses the terminal 5000-call
state followed by ten Newton steps as its center.
The radius is 1.01 times the maximum distance of the three original
154-call outputs from this center, with floor $10^{-9}$.
This frozen rule is a post-hoc certificate construction.
It uses observations, model and states, never truth patches.
Each ball is checked with the reference-component and pairwise
whole-ball covariance bounds.
All 3072 observations satisfy the general sufficient test.
The original uniform choice $\rho=.45$ passes for
768, 765, 765 and 764 observations on the four photos.
For the other ten we propose sample-specific upper bounds by
rounding the floating-point pairwise bound upward to the next
$.001$ and adding $.001$.
This only selects proof constants, and the interval inequalities
must still verify them.

Using 60-digit outward-rounded interval arithmetic, all 3072
LDL curvature, center residual, coverage and step checks pass.
The maximum verified $\rho$ per photo is $.45,.497,.511,.798$,
giving worst factors below
$.999429,.999478,.999493,.999791$, respectively.
The ten exceptions to $\rho=.45$ are therefore not evidence of
noncontraction.
All three original protocols have the same dominant prior component
at 154, 1000 and 5000 calls on every observation.
The interval ball certificate supplies the stronger statement
about subsequent exact-map trajectories.



Across the three new photos, intermediate/terminal mean squared
output distance decreases from $4.38157\cdot10^{-5}$ at 154 calls
to $1.89053\cdot10^{-9}$ at 5000.
Its RMS ratio is approximately $.00657$.
No sampled floating-point distance exceeds its computed bound.
This observation does not extend the interval theorem to all
floating-point trajectories.
The certificate treats exported binary64 coefficients and
154-call states as exact inputs, leaves earlier propagation and
future rounding open, and gives no probability guarantee for
additional observations.

The script \path{experiments/code/patch_independent_validation.py}
exports the frozen plan with \texttt{--prepare} and runs each photo
with \texttt{--photo} and \texttt{--data}.
The interval verifier is
\path{experiments/code/patch_terminal_interval.py}
with \texttt{--adaptive-rho}.
The per-photo risk, seed differences, proof constants and failures
are retained in the corresponding result directories.
Photo provenance is recorded in
\path{experiments/results/patch_independent_data/provenance.md}.
All comparisons and attribution controls remain available.
The positive claim of an arbitrary-budget restricted-controller
risk lower bound is still proved only for the structured family.
This real-photo experiment validates terminal history loss for a
finite batch and has not established positive photo-patch risk
separation.

\section{A Switching Counterexample on a Development Patch}
\label{sec:raw_patch_switching}

This experiment uses patch 37 from a spatial development split and a
separately fitted eight-component isotropic mixture on raw $4\times4$
patches ($\tau^2\simeq .00442835$). The patch was selected by persistent
terminal output separation under the initial schedules. The substitute
schedules were developed on its original observation $y_*$ using distance
to the intermediate schedule's terminal branch. All policies use
the same periodic blur, $\lambda=.1$, $\gamma=.25$, $x_0=A^\top y$,
and 5000 calls, with $\sigma_L=.02$ after the listed prefix.

\begin{table}[tb]
\centering
\small
\begin{tabular}{@{}lcr@{}}
\toprule
Prefix & Ball, entry by & Mean error\\
\midrule
None & High, 2000 & .08413\\
$.1\times25$ & Low, 2000 & .07460\\
$.023\times500$ & Low, 2075 & .07512\\
$(.022\times120,\ 1\times5)\times8$ & Low, 2775 & .07547\\
$.022\times1000$ & High, 2025 & --\\
$1\times1000$ & High, 2875 & --\\
\bottomrule
\end{tabular}
\caption{Prefixes are scale $\times$ calls. High and low denote
the two terminal balls by their squared error to the truth patch.
Entry calls are certified on the original observation neighborhood.
Mean errors use 512 new noise draws and 5000 calls.}
\label{tab:raw_patch_switching}
\end{table}

\paragraph{Observation neighborhood.}
The entry bounds hold for exact trajectories from $x_0=A^\top y$ for
all $y$ with $\|A^\top(y-y_*)\|_2\le 2\cdot10^{-6}$. At each call,
pairwise log-odds bounds control the posterior covariance on the
current state ball. A shifted version of the covariance comparison in
Section~\ref{sec:terminal_history} bounds the update's Jacobian norm
there by $q_k$. If the current radius is $R_k$ and the interval image
of its center has error radius $\epsilon_k$, the next radius is
bounded by $q_kR_k+\epsilon_k+2\gamma\cdot10^{-6}$. We verify the
inequalities with 60-digit outward rounding. The two terminal balls
are disjoint and invariant, with contraction factors below $.997933$
and $.998941$. Their squared errors differ by at least $.08646$
throughout the balls.

The two fixed-scale prefixes and the switching prefix all use 1000
calls. Switching enters the same ball as the intermediate schedule,
giving a successful policy in $[.02,.022]\cup[.3,\infty)$ on this
example. With the narrower low-only interval $[.02,.020001]$,
all switching policies enter the higher-error ball by call 2000
and remain there.

\paragraph{Independent noise.}
We froze the four schedules in the first four rows before drawing
two batches of 256 independent Gaussian noises of standard deviation
$.01$. The middle and the two substitutes reduce the empirical
5000-call error by $11.33\%$, $10.70\%$, and $10.29\%$ against
terminal-only iteration. These estimates are conditional on the
selected truth and prior. The neighborhood certificate and the
noise experiment address local reachability and fixed-patch error,
respectively. The four-photo frozen validation is reported separately
in Section~\ref{sec:independent_patch}.

The interval runs and noise outputs are recorded in
\path{experiments/results/raw_patch_neighborhood/} and
\path{experiments/results/raw_patch_substitute_validation/}.

\section{A uniform non-alignment bound}\label{sec:rotation_proof}
We prove the main-text angle extension over the template and mass
box in~\eqref{eq:recovery_box} at $\alpha=.001$.
Keep the coefficient bounds~\eqref{eq:ind_coeff}. Let $f=.020$,
$B_0=\operatorname{diag}(\alpha,1)$, and
$\Delta B=B_\varphi-B_0$. For $|\varphi|\le f$ we have
$0\le\Delta B_{11}=-\Delta B_{22}\le f^2$ and
$|\Delta B_{12}|\le f$. Use $e$ and $E$ as in~\eqref{eq:ind_event},
with the rotated matrix. At fixed $e$, the rotated map differs
from the aligned map with $b_0=B_0\mu_z+e$ by
$-\gamma\Delta B(x-\mu_z)$.
The rotated map is nondecreasing in each output's own input
coordinate by the same covariance derivative as in S5, since
$B_{\varphi,jj}\le1$. Thus each face estimate extends across its
rectangle for the corresponding output coordinate.

\paragraph{Bounded common barrier.}
Let $W=[-.5,.3]\times[-.027,.51]$ for class 1.
Write $a_\pm,d_\pm,r_\pm,h_\pm$ for the endpoints of the
respective intervals in~\eqref{eq:recovery_box}.
For any $x=(s,t)\in W$, let $R=w_1/w_0$ and $H=w_2/w_0$.
The third-template half squared-distance gap is at least
\[
 Q_{\min}
 =\tfrac12(a_-^2+d_-^2-a_+-2d_+(.51))=.9002 .
\]
For $v\le.17$ this gives $H\le h_+e^{-Q_{\min}/.17}$,
so $w_2<.01$ and $\bar\mu_{v,1}>-.30$.
For $v\ge2.5$, split the range into $[2.5,4]$, $[4,8]$, and
$[8,\infty)$, each with endpoints $l,u$. On each piece,
\[
 R\ge R_* =r_-e^{-1/l},\qquad H\le H^*=h_+e^{-Q_{\min}/u}.
\]
Here $e^{-Q_{\min}/\infty}=1$. On all three pieces,
$H^*/(1+R_*+H^*)<.481$ and
$(R_*-a_+H^*)/(1+R_*+H^*)>-.30$.
Thus throughout $W$ the high and low ranges both satisfy
$w_2<.481$ and $\bar\mu_{v,1}>-.30$.

On the right face $s=.3$ at low scales, dropping $H$ yields
$\bar\mu_{v,1}\le L(.17)<.287729$.
For high scales split $t$ into $[-.027,.15]$ and $[.15,.51]$.
On a piece $[\ell,u]$ the gap is at most
\[
 Q^*=\tfrac12(a_+^2+d_+^2+.6a_+-2d_*\ell),
 \qquad
 d_*=\begin{cases}d_+&\ell<0,\\d_-&\ell\ge0.\end{cases}
\]
Both upper bounds are positive, so
$H\ge h_-e^{-Q^*/2.5}$ and $R\le r_+$.
Substituting these bounds into $(R-aH)/(1+R+H)$ gives
$\bar\mu_{v,1}<.29300$ and $\bar\mu_{v,1}<.25984$ on the
two pieces, respectively. Denote the respective upper bounds by $M$. With $U=\max(|\ell|,|u|)$
on each high-scale piece and $U=.51$ at low scales, the
right-face drift is at most
\[
 \gamma[.00095+.999(M-.3)+.7f^2+fU].
\]
The three upper bounds (low, first high, second high) are below
$-.00020$, $-.00069$, and $-.0071$, respectively.
The left-face drift is at least
$\gamma[.0015-.00025+.999(.5-.30)-f(.51)]>.047$.
The bottom-face drift is at least
$\gamma[-.005+1.999(.027)-1.5f-.027f^2]>.0047$.
The top-face drift is at most
$\gamma[.005-.51+1.002(2.02(.481)-.51)+1.5f+.51f^2]<-.0030$.
For class 1 the initial coordinates on $E$ satisfy
$.00075\le b_1\le.00165$ and $|b_2|\le.025$.
The four margins and the initialization prove that all
restricted trajectories stay in $W$.

\paragraph{Entry under the preset.}
The class-0 and class-1 backprojections retain their original
decision thresholds. For class 2,
$\|\Delta B\mu_2\|_\infty\le(f+f^2)2.02$ implies
$|b_1|<.043$ and $1.933<b_2<2.067$. The classifier therefore
identifies class 2 and its backprojection belongs to
$G_{2,\varphi}=[-1.12,.05]\times[1.73,2.27]$.
Class 0 starts in $G_0$. Both classes use only $v_L$.

For class 1, write $s_-=-.01$, $\xi=.53$, $s_+=1.1$,
$T=.06$, and $v_-=.69,v_+=.71$.
The channel and prepared strip share $-.027\le t\le T$ and have
horizontal intervals $[s_-,\xi]$ and $[\xi,s_+]$, respectively.
On their union,
\begin{align*}
 R_{\rm lo}&=r_-\exp((s_--.5)/v_-),\\
 H_{\rm hi}&=h_+\exp[-(a_-^2+d_-^2+2a_+s_--2d_+T)/(2v_+)]
\end{align*}
bound the two odds from below and above, respectively. Thus
$w_2\le H_{\rm hi}/(1+R_{\rm lo}+H_{\rm hi})<.04508$,
and $\partial_s\bar\mu_{v,1}\le(1/4+a_+(a_++2)w_2)/v_-<.564$.
At $s=\xi$, the lower bound on $R$ becomes
$r_-e^{(\xi-.5)/v_+}$ and the upper bound on $H$ is obtained
by replacing $2a_+s_-$ by $2a_-\xi$. These give a surplus
$\bar\mu_{v,1}(\xi,t)-\xi>.01563$.
The horizontal and vertical rotation corrections on the union
are bounded by $\gamma(1.01 f^2+Tf)<.000402$ and
$\gamma(1.01f+Tf^2)=.005056$, respectively.
Using the $w_2$ bound, the top-face displacement is below
$-.0009$, and the bottom-face displacement exceeds $.007$.
The aligned data force is nonnegative up to $\xi$.
The lower comparison from $\underline s_0=.00075$ is
\[
 \underline s_{k+1}
 =\underline s_k+\gamma(.999)
       [.01563+.436(\xi-\underline s_k)]-.000402 .
\]
This yields $\underline s_{25}>.5306$. The same positive drift
on $s=\xi$ keeps the left face of the prepared strip invariant.
Since $\bar\mu_{v,1}\le1$, the overshoot is below $.649$ and
the right face $s=1.1$ is inward even after adding $.000402$.
The prepared strip lies in $G_1$, so the 25-call switch is valid
for every intermediate history in $[.69,.71]$.
The numerical face and comparison bounds above follow by
outward-rounded evaluation of the displayed scalar expressions.

\paragraph{Actual terminal dynamics and risk.}
The enlargement to $G_{2,\varphi}$ preserves the gap bound $.06$:
on that rectangle the gap against $\mu_0$ is at least
$2(1.98)(1.73)-2(1.02)(.05)-1.02^2-2.02^2>1.62$,
and the gap against $\mu_1$ adds $1-2s\ge.9$.
The other rectangles retain their original gaps.
All terminal rectangles have $\|x-\mu_z\|_\infty\le1.1$.
The rotation correction is bounded there by
$\gamma(f+f^2)1.1=.00561$, which fits the aligned terminal
face margin $\gamma(.999\cdot.09-.005-1.002\cdot12e^{-30})>.021$.
Hence the rotated terminal map preserves all three rectangles.

Use the actual center
$m_\varphi=(B_\varphi+c_LI)^{-1}(b+c_L\mu_z)$.
The residual recursion uses $I-\gamma(B_\varphi+c_LI)$, whose
norm is at most $.751$, and forcing
$\gamma c_L(\bar\mu_{v_L}-\mu_z)$ of norm at most
$12\gamma c_Le^{-.03/v_L}$. Summing the recursion yields
$O(q^n+e^{-.03/v_L})$ with $q=.751$.
The actual terminal map also contracts on the rectangles by
$\|C_{v_L}\|\le36e^{-.03/v_L}$ and $v_L\le.001$.

Conditional on the class, $e$ has covariance
$\tau^2B_\varphi^2+\nu^2B_\varphi\preceq S^2I$,
so $\Prob(E^c)\le4\exp[-(.00025)^2/(2S^2)]$ uniformly in $\varphi$.
The affine-center risk is unchanged because the two sensing
eigenvalues and isotropic covariances are preserved.
On all observations and at every admissible scale the linear
part has norm at most $.751$ and the forcing is bounded by
a constant times $1+\|b\|$. Iterates and centers have
budget-uniform fourth moments for $S\le1$.
Cauchy--Schwarz controls the squared residual on the
exceptional event. Taking $\lambda=1$, $v_L=S$, and the
terminal budget of~\eqref{eq:fixed_geom} yields $O(S^2)$.
On $E$ for class 1 the restricted output still satisfies $s\le.3$,
so the same directional lower bound gives the limiting $.147$.

\paragraph{Reproduction.}
From the artifact root, run
\path{python experiments/code/rotation_certificate.py}.
The only nonstandard dependency is \texttt{mpmath} (1.3.0).
The expected output is \texttt{PASS: 57 strict scalar bounds at
45 digits}. The optional flag \path{--out bounds.json} records
both endpoints of each comparison. The check uses outward-rounded
arithmetic on the scalar reductions above, without trajectory
sampling. The argument relies on the stated fixed gain and
gives the risk rate in~\eqref{eq:fixed_geom}; the calibrated
oracle excess estimate has its own hypotheses.

\section{Scale selection before trajectory evaluation}
\label{sec:selection_diagnostic}
We use $d\in\{1.6,1.8,2,2.2,2.4\}$ and
$h\in\{1.2,1.5,1.78,2.1,2.4\}$, with $a=1$, $r=1.29$,
$\alpha=.001$, $\tau=\nu=\alpha^2$, $v_L=\tau^2+\alpha^2$,
and $\lambda=1+\alpha^2$. The candidate grid is prespecified
relative to this evaluation. We use the S4 tolerances on $E$
and an allowance of $10^{-4}$ for the selection inequalities.
The trajectory simulations use the original RED update in double
precision.

\paragraph{Selection.}
We take the largest $V_-$ in $\{.08,.09,\dots,.30\}$ passing
the low-scale bound of S5 and the smallest $V_+$ in
$\{1.5,1.6,\dots,8\}$ passing its high-scale bound. Both bounds
cover the full restricted ranges $[v_L,V_-]\cup[V_+,\infty)$.
For the channel we take the first passing width-$.02$ band
in increasing order of its center in $\{.40,.42,\ldots,1.10\}$.
We check class-1 progress, transverse invariance, overshoot,
and invariance of the prepared strip through call 25.
The progress comparison starts from $\alpha-\alpha/4$.
Classes 0 and 2 start in their terminal rectangles and use only
$v_L$ under the observation-based preset. We take
$G_2=[-1.1,.01]\times[d-.27,d+.27]$ and use the main-text
$G_0,G_1$. We check the classifier, initialization, positive
template gaps at all corners, and every terminal face inequality.
The bound calculations use 45-digit outward-rounded interval
arithmetic. They hold on $E$ for each certified grid model.
The cases $(d,h)=(1.6,2.1)$ and $(1.6,2.4)$ do not pass a
candidate band under these tests.

\paragraph{Budget-based selection.}
The first-pass rule selected $[.39,.41]$ at all 23 passing
models. After that diagnostic we froze a second rule to
study model-dependent selection. It extends the candidate
centers down to $.24$, and among passing bands minimizes the
first call at which the analytic lower comparison reaches
$.53$, breaking ties by the smaller center. This uses the
model inputs and the 25-call cap. The actual experimental
prefix remains 25 calls for all policies. The new rule
certifies the same 23 settings, with selected centers
between $.36$ and $.60$ and upper bounds of 14 to 24 calls.
It is evaluated on a new noise stream. We execute the
shorter certified prefixes in the follow-up below.

\paragraph{Frozen evaluation.}
After writing each plan, we draw 32 independent pairs of
within-component and measurement noise per class and model,
using seed 20261112 for the first-pass rule and seed
20261113 for the budget-based rule. All 2208 observations
in each evaluation lie in $E$. All policies share their
observations, backprojection, step, terminal variance and
106-call budget. We use 25 prefix calls with alternating
restricted cutoffs, 25 geometric variances from $V_+$ to
$V_-$, or a constant platform at the selected center.
The preset assigns the constant platform only when the
observation classifier predicts class 1. Class labels enter
evaluation and prior weighting. The original terminal map
is then used for at least 81 calls. We retain first entry
times, class counts at the common switch, and class and model
squared errors. Model averages describe the certified grid
subset. The other two models remain uncertified.

Both designs yield the same class-entry counts: restricted
switching has zero passing models, both shared geometric and
shared-platform prefixes pass class 1 on 23 models but not all
classes, and the observation preset passes all classes on all 23.
The first-pass mean errors in this order are
$.32154626,.24926068,.24926068$ and $1.5293\cdot10^{-12}$.
Every class-1 draw under the budget-based rule reaches $G_1$
within its certified bound. These are allocation ablations;
the current main-text table instead reports the independent,
budget-matched controls of S29.

The prespecified failure examples both use $(d,h)=(2,1.78)$.
They hold $V_+=2.2$ for 12 calls and then use $V_-=.18$,
with or without one call at the selected center at the start
of the low stage. Both produce zero class-1 terminal entries
in 106 calls. For the budget-based band $[.43,.45]$, one band
call leaves class 1 at $s<.194$ at the switch. The shared
geometric grid reaches the correct class-1 rectangle on this
fixture at call 19 without using a variance in $[.43,.45]$.
The experiment checks a sufficient prescription, and nearby
successful schedules can use other intermediate variances.

\subsection{Independent platform tuning and realized savings}
We freeze an additional plan on the same 23 certified grid
models. The shared and observation-gated platforms have
identical model inputs. Each searches
$v\in\{.20,.21,\ldots,1.10\}$ and prefix lengths $m=1,\ldots,25$
independently, giving $91\times25=2275$ candidates per
model and policy. For each candidate, run RED from all
three noiseless backprojections $B\mu_z$ with 106 total calls.
The gated policy uses the same classifier as the certified
preset. A candidate passes when the maximum final squared
distance to the corresponding template is at most $10^{-8}$.
Select the shortest passing prefix and the lower median
of its passing variances. Every search has a passing
candidate. The shared choices have $v\in[.41,.85]$ and
the gated choices $v\in[.47,.81]$. Both choose 11--14 calls.

On the noiseless grid the shared search passes 9--25
of the 2275 pairs per model, versus 428--1142 for gating.
At fixed length 25, the counts are 0--2 of 91 for sharing
and 48--86 of 91 for gating. These are discrete measures
of sensitivity to the tuning choices on the model proxies.
Only the frozen choices are carried to the noisy evaluation,
apart from the explicit nominal comparison at $v=.25$.

With seed 20261114 we draw new independent pairs of component
and measurement noise, 64 per class and model, except for the
nominal $(2,1.78)$ model where we draw 512 per class.
All 5760 observations lie in $E$. The certified preset
uses its selected center and analytic cap $m_{\rm cert}$.
All main comparisons use the same 106 calls, observations,
initialization, exact RED formula and terminal variance.
The shorter-budget variant stops at $m_{\rm cert}+81$ on the
same observations. The class labels are used for evaluation
and prior weighting; the sample-specific gate sees only $b$.

\begin{table}[t]
\centering\small
\setlength{\tabcolsep}{2.8pt}
\begin{tabular}{@{}lrrrr@{}}
\toprule
Prefix policy & $G_1$ & All $G_z$ & Mode & Error\\
\midrule
Shared, same center & 23 & 0 & 0 & $.249$\\
Preset, 25-call cap & 23 & 23 & 23 & $1.47\cdot10^{-12}$\\
Preset, certified cap & 23 & 23 & 23 & $1.47\cdot10^{-12}$\\
Shared, tuned & 0 & 0 & 23 & $1.47\cdot10^{-12}$\\
Gated, tuned & 0 & 0 & 23 & $1.47\cdot10^{-12}$\\
Preset, shorter budget & 23 & 23 & 23 & $1.47\cdot10^{-12}$\\
\bottomrule
\end{tabular}
\caption{Frozen follow-up on 23 models. Columns $G_1$ and All
$G_z$ count models whose indicated draws all lie in the named
rectangles after the prefix window. Mode counts models with every
draw's final nearest template correct. Error is prior-weighted
squared error averaged across models.
All rows use 106 total calls except the last (95--105 calls).}
\label{tab:tuned_platforms}
\end{table}

The sufficient rectangles are stricter than the final nearest
template test. On the nominal model, the shared 25-call platform
at total variance $.25$ switches with horizontal coordinates
near $.4961$ and $.5018$ for classes 0 and 1, outside $G_0$ and
$G_1$. It still yields $1536/1536$ correct nearest templates and
prior-weighted squared error $1.5020\cdot10^{-12}$, while the
shared $.44$ prefix yields $1024/1536$ and $.24521$. The tuned
shared choice $(v,m)=(.61,12)$ also yields $1536/1536$ and
$1.5020\cdot10^{-12}$. The shared outcomes are numerical
checks on the selected grid, without an interval guarantee.

For the certified policy, the mean prefix cap drops from 25
to 17.22. Averaging the actual intermediate calls first with
the class priors and then over models gives 8.05 and 5.54 calls
before and after this change. The shorter run saves a mean of
7.78 total calls at the observed error in Table~\ref{tab:tuned_platforms}.

\paragraph{Reproduction.}
The two phases of
\path{experiments/code/scale_selection_diagnostic.py}
write a plan and evaluate it from a separate command.
The frozen plans and per-model outputs are in
\path{experiments/results/scale_selection} and
\path{experiments/results/scale_selection_budget}.
The follow-up uses \path{experiments/code/scale_selection_followup.py}
with \texttt{plan} and then \texttt{evaluate}, both with
\path{--out experiments/results/scale_selection_followup}
when run in a new artifact tree without the reference outputs.
The first command reads the frozen budget-based plan by default.
The second reports the policy summaries in
Table~\ref{tab:tuned_platforms} and writes the per-model output.

{\small\bibliographystyle{ieeenat_fullname}\bibliography{references}}
\end{document}
